% ============================================================================
%  main.tex --- Every Markov model in the thesis, extracted and specified.
%
%  Source: full_draft_thesis/main.pdf ("Stochastic models of early Yersinia
%  pestis infection", Adam Aldridge, Leeds, September 2026 working draft).
%  Compiled from the chapter sources under full_draft_thesis/chapters/.
%
%  Compile from this directory:  latexmk -pdf main.tex
%
%  Self-contained: no thesis macros or preamble required. Thesis \label{}
%  keys are quoted as stable identifiers next to each result.
% ============================================================================
\documentclass[11pt,a4paper]{article}

\usepackage[a4paper,margin=1in]{geometry}
\usepackage{amsmath,amssymb,amsthm}
\usepackage{booktabs}
\usepackage{longtable}
\usepackage{array}
\usepackage{enumitem}
\usepackage[colorlinks=true,linkcolor=blue!45!black,urlcolor=blue!45!black]{hyperref}
\usepackage{microtype}

\newcommand{\PR}{\mathbb{P}}
\newcommand{\EX}{\mathbb{E}}
\newcommand{\Var}{\operatorname{Var}}
\newcommand{\Cov}{\operatorname{Cov}}
\newcommand{\ee}{\mathrm{e}}
\newcommand{\dd}{\mathrm{d}}
\newcommand{\Ind}{\mathbf{1}}
\newcommand{\Lap}{\mathcal{L}}      % Laplace transform
\newcommand{\Rst}{\mathsf{R}}       % two-type absorbing marker
\newcommand{\Geom}{\operatorname{Geom}}
\newcommand{\Exp}{\operatorname{Exp}}

\setlist[itemize]{leftmargin=1.4em,itemsep=1pt,topsep=2pt}

\title{Every Markov model in the thesis\\[4pt]
\large A complete extraction and specification, chapter by chapter}
\author{Extracted from \textit{Stochastic models of early \textit{Yersinia pestis} infection}\\
(full draft, September 2026)}
\date{September 2026}

\begin{document}
\maketitle

\begin{abstract}
\noindent This document extracts every stochastic model with Markovian
dynamics that is defined or substantially used in the thesis, and specifies
each one in a uniform format: class, state space, transition structure
(rates or probabilities), parameters, initial conditions, principal closed
forms, and location in the thesis (chapter, section, and \texttt{label}
key). Thirty-eight numbered models are catalogued, together with four
general frameworks (discrete- and continuous-time Markov chains,
time-inhomogeneous chains, and coupled ODE--CTMC systems) and a short
account of the solution methods the thesis builds around them. Deterministic
companions (mean-field equations, maps, and limits) are included where they
carry a Markov model's exact first moments or serve as its comparator, and
are flagged as such. The document is self-contained and uses standard
notation; the thesis's own \texttt{label} keys are quoted so every entry can
be traced to its source.
\end{abstract}

\tableofcontents

% ============================================================================
\section{Scope, method, and how to read an entry}
\label{sec:scope}
% ============================================================================

\paragraph{What counts as a model here.} An entry is included if the thesis
endows a stochastic process with an explicit transition structure --- a
transition matrix, a list of infinitesimal transition rates, an offspring
distribution, or jump rules of an individual-based system --- or if it is a
deterministic object that carries the exact first moments of such a process
(the burst-aware renewal system) or serves as its named comparator (the
basic model of viral replication, the logistic map). Deterministic,
non-Markovian entries are flagged \textbf{[D]} in their class line. Purely
analytic machinery with no process attached (the hypergeometric integral
identity of Appendix 2.C, coefficient extraction of Appendix 2.D, the
Koenigs function theory) is recorded inside the model entries it serves,
not as a model of its own.

\paragraph{Source.} All material comes from the stitched draft
\texttt{full\_draft\_thesis/main.pdf}, read from its LaTeX chapter sources
(\texttt{chapters/ch01}--\texttt{ch10}). Because the thesis cross-references
by \texttt{label} key rather than printed number, each entry cites the
stable keys (e.g.\ \texttt{bdc:thm:I}, \texttt{dist:thm:qsd}); printed
numbers are avoided except where the conclusion chapter fixes them.

\paragraph{Entry format.} Every model entry follows the same sheet:
\begin{itemize}
  \item \textbf{Class} --- DTMC, CTMC, branching process, PDMP, automaton,
    or \textbf{[D]} deterministic companion;
  \item \textbf{State space} --- exactly as the thesis gives it;
  \item \textbf{Dynamics} --- the defining transition probabilities or
    rates, in the thesis's own equations (house macros normalised to
    $\PR,\EX,\Var$);
  \item \textbf{Parameters} --- symbols and meaning;
  \item \textbf{Initial conditions};
  \item \textbf{Principal results} --- the key closed forms and laws
    derived for the model, with thesis \texttt{label} keys;
  \item \textbf{Location and connections} --- where it lives and how it
    relates to the other entries.
\end{itemize}

\paragraph{Notation.} $\PR,\EX,\Var,\Cov$ have their usual meaning;
$\Exp(\theta)$ is an exponential holding time; $\Lap f(r)=\int_0^\infty
\ee^{-r\alpha}f(\alpha)\dd\alpha$ is a Laplace transform; $\Ind\{\cdot\}$ an
indicator. For the central BDC process the thesis's derived constants are
used throughout: with
\[
\eta=\frac{\lambda+\mu+\delta}{2\lambda},\qquad
a=\eta+\sqrt{\eta^{2}-\tfrac{\mu}{\lambda}},\qquad
b=\eta-\sqrt{\eta^{2}-\tfrac{\mu}{\lambda}},
\]
\[
A:=a-1,\qquad B:=1-b,\qquad \theta:=\lambda(a-b),\qquad
\kappa:=1+\frac{\delta}{2\lambda},\qquad w:=\ee^{\theta t},
\]
so that $b<1<a$, $ab=\mu/\lambda$, $A+B=a-b$ and $AB=\delta/\lambda$
(\texttt{bdc:lem:roots}). The symbol $L=a(1-b)=a-\mu/\lambda$ is the
flooding parameter of Chapter 6. $\Rst$ denotes the two-type chapter's
marker for the absorbing catastrophe state.

\paragraph{The spine.} One process organises the collection: the linear
birth--death--catastrophe process (M6) is defined in Chapter 2, solved in
Chapter 4, given full distribution theory in Chapter 5, lifted to the
population level in Chapter 6, extended to two types in Chapter 7, and
passed through selection functionals in Chapter 8. Around it sit the
discrete-time Galton--Watson branch (M2), the classical chains of the
mathematical introduction (M1, M3--M5), the two-compartment absorption
family (M13--M15), and the non-constant-rate models of Chapter 9.

% ============================================================================
\section{Master inventory}
% ============================================================================

{\small
\begin{longtable}{@{}p{0.9cm}p{4.3cm}p{2.6cm}p{2.2cm}p{4.6cm}@{}}
\toprule
\textbf{ID} & \textbf{Model} & \textbf{Class} & \textbf{Chapter} &
\textbf{Defining feature / role} \\
\midrule
\endhead
F1 & Discrete-time Markov chain & framework & 2 & $P_{ij}$, Chapman--Kolmogorov \\
F2 & Continuous-time Markov chain & framework & 2 & generator $q_{ij}$, exponential clocks \\
F3 & Time-inhomogeneous CTMC & framework & 2 & $Q(t)$, Kolmogorov equations \\
F4 & Coupled ODE--CTMC (PDMP) & framework & 2 & joint generator, three coupling types \\
\midrule
M1 & Random walk / gambler's ruin & DTMC & 2 & $P_{i,i\pm1}$; ruin probabilities \\
M2 & Galton--Watson (binary) & DTMC branching & 1--3 & offspring $p_2=p$; $A(p)$, Koenigs \\
M3 & Poisson process & CTMC & 2 & $q_{n,n+1}=\alpha$ \\
M4 & Yule process & CTMC & 1, 2 & $q_{n,n+1}=\lambda n$ \\
M5 & Linear birth--death process & CTMC & 1, 2 & $\lambda n,\mu n$; Kendall solution \\
M6 & Birth--death--catastrophe (BDC) & CTMC & 2, 4, 5 & $\lambda n,\mu n,\delta n\!\to\!H$; the spine \\
M7 & Discrete-generation catastrophe & DTMC variant & 2 & per-generation hazard $\chi$; Jensen gap \\
M8 & Inhomogeneous linear BD & CTMC, $Q(t)$ & 2 & $\lambda(t),\mu(t)$; Riccati \\
M9 & Seasonal Lotka--Volterra & 2-pop.\ CTMC & 2 & periodic prey birth $\mathcal S^+(t)$ \\
M10 & Logistic speciation & nonhom.\ Yule & 2, 9 & $\beta(t)=\sigma(K-N(t))$ \\
M11 & Gated release & PDMP & 2 & gate opens at rate $\eta A_t$ \\
M12 & Surging gated release & PDMP & 2 & gate opens at rate $\nu$ \\
M13 & Absorption-only & bivariate CTMC & 2 & uptake $\alpha$ only \\
M14 & Absorption--death & bivariate CTMC & 2 & uptake $\alpha$, death $\mu$ \\
M15 & Absorption--birth--death & bivariate CTMC & 2 & $+$ interior BD; novel pgf \\
\midrule
M16 & Chained immediate transfer ($\mu=0$) & chained CTMC & 5, 8 & burst founds next cell \\
M17 & Budding comparator cell & two-stream Poisson & 5, 6 & release $p$, death $d_I$ \\
\midrule
M18 & Basic model of viral replication & \textbf{[D]} ODE & 1, 6 & $\dot I=\gamma TV-d_II$, $\dot V=pI-cV$ \\
M19 & Within-host population process & CTMC population & 6 & BDC cells $+$ virion infect/clear \\
M20 & Burst-aware renewal system & \textbf{[D]} exact means & 6 & $I=I_0\hat I+(i*\hat I)$, $V'=I_0g+(i*g)-cV$ \\
M21 & Effective-parameter projection & \textbf{[D]} map & 6 & $p_{\rm eff}(r)$, $d_{I,\rm eff}(r)$ \\
M22 & Generation kernel (Euler--Lotka) & \textbf{[D]} functional & 6 & $\mathcal A(\alpha)$; $R_0=\gamma TV_\infty/c$ \\
M23 & Establishment branching process & GW offspring & 6 & burst vs bud offspring; flooding $L$ \\
M24 & Invasion-speed comparison & \textbf{[D]} & 6 & $r_{\rm bud}>r_{\rm burst}$; Laplace order \\
M25 & Reset process & CTMC $+$ death & 6 & dumps to $0$ at rate $\delta n$, cell dies $d_I$ \\
M26 & Partial-release stock model & \textbf{[D]} mean-field & 6 & boolean thinning $\varphi$; stock $Q$ \\
M27 & HIV skeleton variant & specification & 6 & Allee incidence, Erlang eclipse \\
\midrule
M28 & Two-type catastrophe process & 2-type CTMC & 7 & $\lambda_i,\mu_i,\nu$; $\delta_1x+\delta_2y\to\Rst$ \\
\midrule
M29 & Replicator--suppressor & nonlinear CTMC & 8 & birth $\lambda i$, kill $\varsigma ij$ \\
M30 & Playing-field particle model & spatial simulation & 8 & Brownian; killers used up \\
M31 & Fixed-removal ratio model & functional on M6 & 8 & $\mathcal R(\vartheta)=V_\infty a^{-\vartheta}$ \\
\midrule
M32 & Pimentel competition chain & path-dep.\ DTMC & 9 & moving unstable critical point \\
M33 & Pimentel$+$ & specification & 9 & fitness $+$ species BD coupling \\
M34 & Rise, run, ruin, rejuvenation & CTMC $+$ ODE & 9 & catastrophe $\delta X_t/V(t)$ \\
M35 & Multiplicative dissipation & stochastic CA & 9 & torus $256^2$; fragmentation $\delta X_i/V_i$ \\
M36 & Birth--plague ($+$ vitality) & two-type CTMC & 9 & bilinear infection $\chi X_tY_t$ \\
M37 & Public good accumulation & CTMC $+$ ODE & 9 & $\lambda(t)=\lambda_0+\alpha P(t)$ \\
\midrule
M38 & Quadrant argument & \textbf{[D]} comparison & 10 & threshold $p_N^*=(\mu_E-\mu_M)/(\mu_N-\mu_M)$ \\
\bottomrule
\end{longtable}
}

% ============================================================================
\section{Frameworks (Chapter 2)}
% ============================================================================

The mathematical introduction defines four general containers before any
instance; every catalogue entry is an instance of one of them.

\subsection*{F1 --- Discrete-time Markov chain (\texttt{m:sec:dtmc})}
Time-homogeneous DTMC on a countable $\mathcal S$ with
\[
\PR{X_{n+1}=j\mid X_n=i,X_{n-1},\ldots,X_0}
=\PR{X_{n+1}=j\mid X_n=i}=:P_{ij}
\qquad\text{(\texttt{m:eq:dtmarkov})}
\]
and the Chapman--Kolmogorov composition
$(P^{n+m})_{ij}=\sum_k (P^n)_{ik}(P^m)_{kj}$
(\texttt{m:eq:chapmankolmogorov}). Absorbing states have $P_{ii}=1$.

\subsection*{F2 --- Continuous-time Markov chain (\texttt{m:sec:ctmc})}
Generator
\[
q_{ij}=\lim_{\Delta t\downarrow0}\frac{\PR{X_{t+\Delta t}=j\mid X_t=i}}{\Delta t}\quad(j\ne i),
\qquad q_{ii}=-\sum_{j\ne i}q_{ij}=:-q_i
\qquad\text{(\texttt{m:eq:generator})}
\]
so $\PR{X_{t+\Delta t}=j\mid X_t=i}=q_{ij}\Delta t+o(\Delta t)$; state $i$
is absorbing iff $q_i=0$. Holding times are exponential
(\texttt{m:eq:expdef}); the competing-clocks proposition
(\texttt{m:prop:race}) gives $\min_k T_k\sim\Exp(\sum_k\theta_k)$ with
$\PR{T_k\text{ wins}}=\theta_k/\sum\theta_k$, and its consequence is
Gillespie's direct method (\texttt{m:rem:gillespie}).

\subsection*{F3 --- Time-inhomogeneous CTMC (\texttt{m:sec:inhomogeneous})}
Two-time kernels $p_{ij}(s,t)$ with matrix Kolmogorov equations
\[
\partial_t P(s,t)=P(s,t)Q(t),\qquad \partial_s P(s,t)=-Q(s)P(s,t)
\qquad\text{(\texttt{m:eq:inhomogkolmogorov})}.
\]

\subsection*{F4 --- Coupled ODE--CTMC systems / PDMPs (\texttt{m:sec:coupled})}
Continuous state $\mathbf y\in\mathbb R^d$, discrete mode $M_t\in\mathcal S$,
joint generator
\[
(\mathcal Lf)(\mathbf y,i)
=F(\mathbf y,i)\cdot\nabla_{\mathbf y}f(\mathbf y,i)
+\sum_{j\ne i}q_{ij}(\mathbf y)\bigl[f(\mathbf y,j)-f(\mathbf y,i)\bigr]
\qquad\text{(\texttt{m:eq:coupledgenerator})}
\]
with three coupling types: (i) autonomous chain switches the vector field;
(ii) a prescribed trajectory drives the rates,
$q_{ij}(t)=q_{ij}(\mathbf y(t))$; (iii) two-way coupling. Holding-time law
for simulation: $\PR{T>s\mid\mathbf y_t,M_t=i}
=\exp\{-\int_t^s q_i(\mathbf y_u)\dd u\}$ (\texttt{m:eq:coupledholding}).

% ============================================================================
\section{Catalogue, part I --- the mathematical introduction (Chapters 1--2)}
% ============================================================================

Chapter 1 (\S1.4) introduces the binary Galton--Watson process, the Yule
process, the standard birth--death process and the birth--death--catastrophe
process informally through their infinitesimal rates; the formal definitions
are those of Chapter 2 reproduced below, and \S1.6 introduces the
deterministic BMVR (M18). The entries below give the formal specifications.

% ---------------------------------------------------------------
\subsection{M1 --- Simple random walk and gambler's ruin}
\label{sec:M1}

\paragraph{Class.} DTMC on the integers.

\paragraph{State space.} $\mathbb Z$; for the ruin problem
$\{0,1,\ldots,M\}$ with $0,M$ absorbing.

\paragraph{Dynamics.}
\[
P_{i,i+1}=q,\qquad P_{i,i-1}=1-q
\qquad\text{(\texttt{m:eq:rwrates})}.
\]

\paragraph{Parameters.} $q\in(0,1)$, symmetric at $q=\tfrac12$; initial
state $X_0$ deterministic.

\paragraph{Initial conditions.} $X_0=i$.

\paragraph{Principal results.} $\EX{X_n}=X_0+n(2q-1)$, $\Var{X_n}=4nq(1-q)$
(\texttt{m:eq:rwmoments}); with $R\sim\operatorname{Bin}(n,q)$,
$X_n=X_0+2R-n$ (\texttt{m:eq:rwlaw}). Hitting probability of $M$ before $0$,
via the first-step relation $h_i=qh_{i+1}+(1-q)h_{i-1}$, $h_0=0$, $h_M=1$
(\texttt{m:eq:gamblersruin}):
\[
h_i=\frac{1-\theta^i}{1-\theta^M}\ \ (\theta=\tfrac{1-q}{q}\ne1),
\qquad h_i=\frac{i}{M}\ \ (q=\tfrac12).
\]

\paragraph{Location and connections.} Chapter 2, \S2.2.1
(\texttt{m:sec:randomwalk}). Elementary instance of the first-step
principle (\texttt{m:prop:firststep}); the embedded jump chain of the linear
birth--death process M5 with $q=\lambda/(\lambda+\mu)$
(\texttt{m:eq:embedded}).

% ---------------------------------------------------------------
\subsection{M2 --- Galton--Watson branching process (binary specialisation)}
\label{sec:M2}

\paragraph{Class.} Discrete-time Galton--Watson branching process; ``the
principal discrete-time process in this thesis''.

\paragraph{State space.} $Z_n\in\{0,1,2,\ldots\}$, population census;
$0$ absorbing.

\paragraph{Dynamics.} $Z_0=1$; each of the $Z_n$ individuals independently
leaves an iid number of offspring $L$ with $\PR{L=k}=p_k$,
\[
Z_{n+1}=\sum_{i=1}^{Z_n}L_i
\qquad\text{(\texttt{m:eq:gwrecursion})}.
\]
Binary specialisation used throughout (\texttt{m:eq:binarylaw}): replicate
with probability $p$, die with probability $1-p$,
\[
p_2=p,\qquad p_0=1-p,\qquad p_k=0\ (k\notin\{0,2\}),\qquad
\phi(z)=\EX{z^L}=pz^2+(1-p).
\]

\paragraph{Parameters.} $p\in(0,1)$; mean offspring $m=\EX{L}=2p$;
regimes subcritical/critical/supercritical according as $2p\lessgtr1$.
Chapter 3 writes $\varepsilon:=1-2p$ for the distance from criticality.

\paragraph{Initial conditions.} $Z_0=1$ (single founder); cohort version
$Z_0=k$ via independence.

\paragraph{Principal results.}
\begin{itemize}
\item \emph{PGF iteration} (\texttt{m:eq:gwcomposition}): $\phi_{n+1}=
  \phi_n\circ\phi$, $\phi_0(z)=z$, hence $\EX{Z_n}=(2p)^n$.
\item \emph{Extinction/survival recursions} (\texttt{m:eq:unrec},
  \texttt{eq:SnIter}): with $u_n=\PR{Z_n=0}$, $S_n=1-u_n$,
  \[
  u_{n+1}=(1-p)+pu_n^2=\phi(u_n),\qquad
  S_{n+1}=2pS_n-pS_n^2,\qquad S_0=1.
  \]
  Fixed points $0$ and $S_\infty=(2p-1)/p$ (positive only for $p>\tfrac12$;
  extinction certain otherwise).
\item \emph{The constant $A(p)$} (\texttt{m:eq:Apdef},
  \texttt{eq:Ap-def-main}, \texttt{eq:Sn-asymp}), for $0<p<\tfrac12$:
  \[
  S_n\sim A(p)\,(2p)^n,\qquad
  A(p)=\lim_{n\to\infty}\frac{S_n}{(2p)^n}
  =\frac12\prod_{k=1}^{\infty}\Bigl(1-\frac{S_k}{2}\Bigr),
  \]
  with endpoint $A(0)=\tfrac12$.
\item \emph{Yaglom (quasi-stationary) limits}
  (\texttt{m:eq:qsmean}, \texttt{eq:condvar}):
  \[
  \lim_{n\to\infty}\EX{Z_n\mid Z_n>0}=\frac{1}{A(p)},\qquad
  \lim_{n\to\infty}\Var{Z_n\mid Z_n>0}
  =\frac{1}{A(p)}\Bigl(1+\frac{1}{1-2p}\Bigr)-\frac{1}{A(p)^2}.
  \]
\item \emph{Critical power laws} (\S critical endpoint, Ch.~3):
  $S_n\sim 2/n$, $\PR{T_{\rm extinct}=n}\sim2/n^2$, total progeny tail
  $\PR{\mathcal T=n}\sim n^{-3/2}$.
\item \emph{Near-critical asymptotics}
  (\texttt{sec:nearcrit}, \texttt{eq:nearcrit-profile}):
  \[
  A(p)\sim2\varepsilon=2(1-2p),\qquad
  S_n\simeq\frac{2\varepsilon}{\ee^{\varepsilon n}-1},\qquad
  \EX{Z_n\mid Z_n>0}\sim\frac{1}{2\varepsilon},\qquad
  \Var{Z_n\mid Z_n>0}\sim\frac{1}{4\varepsilon^2}.
  \]
\item \emph{Koenigs identification} (Ch.~3 \texttt{def1},
  \texttt{gw:thm:main}): the rescaling $S_n=2w_n$, $r=2p$ turns
  \texttt{eq:SnIter} into the logistic map
  $w_{n+1}=rw_n(1-w_n)$, $w_0=\tfrac12$, $0<r<1$, and
  \[
  A(p)=2\psi_{2p}\bigl(\tfrac12\bigr),
  \qquad \psi(f_r(z))=r\psi(z),\ \psi(0)=0,\ \psi'(0)=1,
  \]
  with $\psi_r(z)=\lim_{n\to\infty}r^{-n}f_r^{\circ n}(z)$. Theorem
  \texttt{gw:thm:main}: $\psi$ is hypertranscendental for every $0<r<1$
  (via Becker--Bergweiler \texttt{thm:bb} and Lemma \texttt{lem:inverse});
  consequently no elementary formula for $A(p)$ exists through this route.
\item \emph{Cohort version} (Appendix 2.A, \texttt{m:eq:Snk}): from $k$
  founders, $S_n^{(k)}=1-u_n^k$; establishment probability
  $1-\bigl((1-p)/p\bigr)^k$ for $p>\tfrac12$, zero at or below criticality.
\item \emph{Computation} (Ch.~3): ratio monitor, truncation bound, stopping
  rule $S_N\le2\delta(1-2p)$, cost $N\approx\varepsilon^{-1}\ln(1+1/\delta)$,
  and the hybrid estimator switching to $A\approx2-4p$ near criticality.
\end{itemize}

\paragraph{Location and connections.} Introduced informally in Chapter 1
\S1.4; formal definition \S2.2.1 (\texttt{m:sec:gw}); analysed in depth in
Chapter 3 (\texttt{gw:ch:galton-watson}), including the formula search
(\texttt{sec:search}) and visual Koenigs geometry. Continuous-time
counterpart constant $A_c$ of M5. Its offspring-pgf fixed point
$s=\phi(s)$ appears in the first-step section (\texttt{m:eq:extinctionfixedpoint}).

% ---------------------------------------------------------------
\subsection{M3 --- Poisson process}

\paragraph{Class.} CTMC counting process.

\paragraph{State space.} $\{0,1,2,\ldots\}$, no absorbing state.

\paragraph{Dynamics.} $q_{n,n+1}=\alpha$ for $n\ge0$
(\texttt{m:eq:poissonrates}).

\paragraph{Parameters.} $\alpha>0$ event rate. Initial $N_0=0$.

\paragraph{Principal results.}
$\PR{N_t=n}=\ee^{-\alpha t}(\alpha t)^n/n!$ with mean and variance $\alpha t$
(\texttt{m:eq:poissonlaw}); thinning (keep each event independently with
probability $\varpi$) gives rate $\varpi\alpha$ --- the basis for simulating
time-dependent rates.

\paragraph{Location and connections.} \S2.2.2 (\texttt{m:sec:poisson}).
Building block of M17 (two independent streams) and of the thinning method
used for M9 and F3-type rates.

% ---------------------------------------------------------------
\subsection{M4 --- Yule process (standard birth process)}

\paragraph{Class.} CTMC, linear pure birth.

\paragraph{State space.} $\{1,2,\ldots\}$ (no death, no absorption below the
initial count).

\paragraph{Dynamics.} $q_{n,n+1}=\lambda n$, $n\ge1$
(\texttt{m:eq:yulerates}).

\paragraph{Parameters.} $\lambda>0$ per-capita birth rate; single founder
initial condition.

\paragraph{Principal results.} Geometric law with exponential parameter
(\texttt{m:eq:yulelaw}):
\[
\PR{Y_t=n}=\ee^{-\lambda t}\bigl(1-\ee^{-\lambda t}\bigr)^{n-1},\qquad
\EX{Y_t}=\ee^{\lambda t},\qquad
\Var{Y_t}=\ee^{\lambda t}\bigl(\ee^{\lambda t}-1\bigr),
\]
holding time in state $n$ is $\Exp(\lambda n)$; coefficient of variation
$\to1$.

\paragraph{Location and connections.} Introduced in Chapter 1 \S1.4 (Yule
1925); formal treatment \S2.2.2 (\texttt{m:sec:yule}). The $\mu=0$ case of
M5; the constant-rate base of M10 (inhomogeneous Yule); its tree appears as
the BEAST Yule prior (Remark \texttt{m:rem:yuleorigin}).

% ---------------------------------------------------------------
\subsection{M5 --- Linear birth--death process}
\label{sec:M5}

\paragraph{Class.} CTMC; the running example for the methods of \S2.3.

\paragraph{State space.} $\{0,1,2,\ldots\}$; $0$ absorbing.

\paragraph{Dynamics.}
\[
q_{n,n+1}=\lambda n,\qquad q_{n,n-1}=\mu n
\qquad\text{(\texttt{m:eq:bdrates})}.
\]

\paragraph{Parameters.} $\lambda$ per-capita birth, $\mu$ per-capita death;
regime set by $\lambda-\mu$.

\paragraph{Initial conditions.} $X_0=N$.

\paragraph{Principal results.}
\begin{itemize}
\item \emph{Master equation} (\texttt{m:eq:bdmaster}):
  $\dot p_n=\lambda(n-1)p_{n-1}+\mu(n+1)p_{n+1}-(\lambda+\mu)np_n$; pgf PDE
  (\texttt{m:eq:bdpde})
  \[
  \partial_tG=(\lambda z^2-(\lambda+\mu)z+\mu)\,\partial_zG
  =(\lambda z-\mu)(z-1)\,\partial_zG,\qquad G(z,0)=z^N.
  \]
\item \emph{Kendall's solution} (\texttt{m:eq:bdsolution}), $\lambda\ne\mu$:
  \[
  G(z,t)=\Bigl[\frac{\mu(1-z)-(\mu-\lambda z)\ee^{-(\lambda-\mu)t}}
  {\lambda(1-z)-(\mu-\lambda z)\ee^{-(\lambda-\mu)t}}\Bigr]^{N}.
  \]
\item \emph{Moments} (\texttt{m:eq:bdmean}, \texttt{m:eq:bdvariance}):
  $\EX{X_t}=N\ee^{(\lambda-\mu)t}$;
  $\Var{X_t}=N\frac{\lambda+\mu}{\lambda-\mu}\ee^{(\lambda-\mu)t}
  (\ee^{(\lambda-\mu)t}-1)$, critical limit $2N\lambda t$.
\item \emph{Extinction} (\texttt{m:eq:p0t}): $p_0(t)=G(0,t)$ in closed form
  (both regimes); ultimate extinction probability $(\mu/\lambda)^N$ when
  $\lambda>\mu$, certain otherwise; at criticality the mean absorption time
  is infinite.
\item \emph{Quasi-stationarity} (\texttt{m:def:qsd},
  \texttt{m:eq:ctsqsmean}): conditional mean limit
  \[
  \lim_{t\to\infty}\EX{X_t\mid X_t>0}=\frac{\mu}{\mu-\lambda}=:\frac{1}{A_c},
  \qquad
  A_c(p)=\frac{1-2p}{1-p}=\frac{2\varepsilon}{1+\varepsilon},
  \quad p=\frac{\lambda}{\lambda+\mu},
  \]
  i.e.\ the continuous-time analogue of $A(p)$ --- a closed form where M2's
  $A(p)$ is only a limit (\texttt{m:eq:twoconstants}).
\item \emph{Embedded chain}: M1 with $q=\lambda/(\lambda+\mu)$
  (\texttt{m:eq:embedded}); barrier-hitting probabilities
  (\texttt{m:eq:bdruin}) as in M1 with $\theta=\mu/\lambda$.
\end{itemize}

\paragraph{Location and connections.} Chapter 1 \S1.4 (informal), \S2.2.2
(\texttt{m:sec:bd}) with the full solution package in \S2.3.2; the
$\delta\downarrow0$ limit of M6 (Appendix \texttt{bdc:app:nocatastrophe});
base process for M8, M9 (marginal), M10 (time change).

% ---------------------------------------------------------------
\subsection{M6 --- Birth--death--catastrophe (BDC) process}
\label{sec:M6}

The central model of the thesis, in the catastrophe convention.

\paragraph{Class.} CTMC on the non-negative integers with a second
absorbing state; equivalently a linear birth--death process killed at a
load-proportional hazard. Non-explosive (Remark \texttt{bdc:rem:assumptions}).

\paragraph{State space.} $\{0,1,2,\ldots\}\cup\{H\}$: intracellular load,
with $0$ = internal extinction (clearance) and $H$ = rupture/lysis; both
absorbing.

\paragraph{Dynamics.} Definition \texttt{m:def:bdc} / \texttt{bdc:def:bdc} /
\texttt{dist:def:bdc}: each replicative unit independently gives birth at
rate $\lambda$, dies at rate $\mu$, and triggers catastrophe at rate
$\delta$; for $n\ge1$
\[
q_{n,n+1}=\lambda n,\qquad q_{n,n-1}=\mu n,\qquad q_{n,H}=\delta n
\qquad\text{(\texttt{bdc:eq:rates})}.
\]
\emph{Augmentation --- released-count process} (\texttt{m:eq:bdcrelease},
\texttt{bdc:eq:release}): with $\tau=\inf\{t:X_t=H\}$,
\[
W_t=\begin{cases}0,&t<\tau,\\ X_{\tau-},&t\ge\tau,\end{cases}
\qquad K:=X_{\tau-}\ \text{(burst size; the left limit is essential)}.
\]
The joint process $(X_t,W_t)$ has the four-channel transition list
\texttt{bdc:eq:joint} (birth, death, catastrophe with $W$ jumping to $n$,
no event).

\paragraph{Parameters.} $\lambda>0$ birth, $\mu\ge0$ death, $\delta>0$
catastrophe rates per unit; derived constants $\eta,a,b,A,B,\theta,\kappa,w$
as in the notation paragraph of \S\ref{sec:scope}. Interpretations: $b$ = probability of internal
clearance before rupture; $1/a$ = ratio of the geometric burst law.
Figure working point $(\lambda,\mu,\delta)=(1,0.2,0.05)$; published
parameter sets (\texttt{dist:tab:twoextremes}): \emph{F.\ tularensis}
$(0.15,0.01,1.5\times10^{-4})\,\mathrm h^{-1}$, \emph{B.\ anthracis}
$(0.64,1.64,0.04)\,\mathrm h^{-1}$.

\paragraph{Initial conditions.} $X_0=1$ unless stated; $W_0=0$;
$I(0)=1$, $D(0)=0$, $\hat I(0)=1$.

\paragraph{Principal results --- fixation functions (Chapter 4).}
With $I(t)=\PR{X_t\ne H}$, $D(t)=\PR{X_t=0}$,
$\hat I(t)=I-D=\PR{X_t\notin\{0,H\}}$ (``productive''):
\[
\dot I=\mu-(\lambda+\mu+\delta)I+\lambda I^2=\lambda(I-a)(I-b),\qquad I(0)=1
\]
(\texttt{bdc:eq:riccati}, \texttt{bdc:eq:riccatifact}), and
\[
I(t)=\frac{aB+bAw}{B+Aw},\qquad
D(t)=\frac{ab(w-1)}{aw-b},\qquad
\hat I(t)=\frac{(a-b)^2w}{(B+Aw)(aw-b)}
\]
(Theorems \texttt{bdc:thm:I}, \texttt{bdc:thm:D}, \texttt{bdc:thm:Ifix});
all three tend to $b,b,0$.

\paragraph{Principal results --- moments and burst statistics (Chapter 4).}
\[
J(t)=\EX{X_t}=\frac{(a-b)^2w}{(B+Aw)^2}=-\frac{\lambda}{\delta}(I-a)(I-b),
\qquad
K(t)=\EX{X_t^2}=\frac{2\lambda}{\delta}(\kappa-I)J,
\]
\[
V(t)=\EX{W_t}=(1-I)\Bigl[1+\frac{\lambda}{\delta}(1-I)\Bigr],\qquad
V_\infty=\frac{aB}{A}=\frac{a(1-b)}{a-1},
\]
(Theorems \texttt{bdc:thm:J}--\texttt{bdc:thm:V}); $\Cov(X_t,W_t)=-JV\le0$;
open moment hierarchy closed by the polynomial property ``for each $p\ge1$
there is a polynomial $Q_p$ of degree $p+1$ with $\EX{X_t^p}=Q_p(I(t))$''
(Proposition \texttt{bdc:prop:polynomial}). Defective rupture-time density
$\varphi(t)=-\dot I=\delta J(t)$ with total mass $1-b$
(\texttt{bdc:eq:bursttime}); conditional burst moments
\[
\EX{K\mid\text{burst}}=\frac{a}{a-1},\qquad
\Var{K\mid\text{burst}}=\frac{a}{(a-1)^2}
\]
(Corollaries \texttt{bdc:cor:burstmean}, \texttt{bdc:cor:burstvar}).

\paragraph{Principal results --- distribution theory (Chapter 5).}
\begin{itemize}
\item \emph{Master equations and pgf} (\texttt{dist:eq:kolmogorov},
  \texttt{dist:eq:catpde}): $\partial_tG=(\lambda z^2-(\lambda+\mu+\delta)z+
  \mu)\,\partial_zG$ with
  \[
  G(z,t)=\frac{a(z-b)-b(z-a)w}{(z-b)-(z-a)w}
  =\frac{ab(1-w)-z(a-bw)}{(b-aw)-z(1-w)}
  \qquad\text{(\texttt{dist:eq:gcat})}
  \]
  (defective: $G(1,t)=I(t)\le1$).
\item \emph{Geometric state law at every time} (Theorem
  \texttt{dist:thm:geometric}): $p_n(t)=p_1(t)P(t)^{n-1}$ for $n\ge1$ with
  sliding ratio
  \[
  P(t)=\frac{1-w}{b-aw},\qquad P(0)=0,\quad P(\infty)=\frac1a .
  \]
\item \emph{Quasi-stationary distribution} (Theorem \texttt{dist:thm:qsd}):
  conditioning on productivity, the limit is geometric,
  \[
  q_n=\Bigl(1-\frac1a\Bigr)\Bigl(\frac1a\Bigr)^{n-1}=(a-1)a^{-n},\qquad n\ge1,
  \]
  with $\langle X\rangle_{\rm QS}=a/(a-1)$, $\Var_{\rm QS}=a/(a-1)^2$; the
  QSD decay rate is exactly $\theta=\lambda(a-b)$ (Proposition
  \texttt{dist:prop:decay}) and the mean productive lifetime is
  $\EX{T_{\rm prod}}=\lambda^{-1}\log\frac{a}{a-1}$ (Proposition
  \texttt{dist:prop:lifetime}).
\item \emph{Burst-size law} (Theorem \texttt{dist:thm:burst}) and its
  identity with the QSD (Theorem \texttt{dist:thm:identity}):
  \[
  \PR{K=k}=\frac{\delta}{\lambda}\,a^{-k}\ (k\ge1),\quad
  \PR{K=0\ \text{(no burst)}}=b;\qquad
  \PR{K=k\mid\text{burst}}=(a-1)a^{-k}=q_k .
  \]
\item \emph{Burst-time statistics} (Chapter 5 \S6--7): peak of $p_n$ at
  $t_n=(\log n)/\theta$; at $\mu=0$,
  $\EX{\tau\mid K=n}=H_n/\theta\simeq(\log n+\gamma)/\theta$ and
  $\PR{\tau<t\mid K=n}=(1-\ee^{-\theta t})^n\to$ Gumbel
  (Proposition \texttt{dist:prop:taugivenk}); for $\mu>0$,
  $\EX{\tau\mid\text{burst}}=\frac{1}{\lambda(1-b)}\log\frac{a-b}{a-1}$
  (Proposition \texttt{dist:prop:tauburst}); size-biased means
  $\EX{K\mid\tau=t}=K(t)/J(t)\to(a+1)/(a-1)$ (the ``patient-cell'' limit).
\end{itemize}

\paragraph{Variants and limits.}
\begin{itemize}
\item \emph{Killing convention} (Chapter 4 \texttt{bdc:sec:conventions};
  Chapter 5 Proposition \texttt{dist:prop:conventions}): catastrophe sends
  $n\to0$ instead of $n\to H$; burst times and sizes unchanged;
  $G_{\rm kill}(z,t)=G_{\rm cat}(z,t)+(1-I(t))$, so $p_n$ ($n\ge1$) and all
  load moments are convention-independent.
\item \emph{$k$ founders (MOI-$k$)} (\texttt{dist:sec:moi}): $G_k=G^k$;
  $I_k=I^k$, $D_k=D^k$, $\hat I_k=I^k-D^k$ (not $\hat I^k$);
  $J_k=kI^{k-1}J$, $K_k=kKI^{k-1}+k(k-1)J^2I^{k-2}$, release kernel
  $g_k=\delta K_k$; lifetime yield $V_\infty^{(k)}$ in closed form, with
  $V_\infty^{(k)}=k+\lambda/\delta$ at $\mu=0$; yield sub-extensive in $k$.
\item \emph{Degenerate limits} (Appendix \texttt{bdc:app:limits}): $\mu=0$
  gives $b=0$, $a=1+\delta/\lambda$, certain rupture, $V_\infty=1+\lambda/
  \delta$; $\delta\downarrow0$ ($\lambda>\mu$) recovers M5 with
  $D(t)\to\frac{\mu}{\lambda}\frac{w-1}{w-\mu/\lambda}$; $\lambda\downarrow0$
  leaves the two-clock race $\dot I=\mu-(\mu+\delta)I$,
  $b\to\mu/(\mu+\delta)$.
\end{itemize}

\paragraph{Location and connections.} Informal rates in Chapter 1 \S1.4;
Definition \texttt{m:def:bdc} (\S2.2.2) with rate diagram
\texttt{m:fig:ratediagram}; core theory Chapter 4 (\texttt{bdc:ch:core},
\S\S2--4 and Appendix A); distribution theory Chapter 5
(\texttt{dist:ch:distribution-theory}); exported to Chapter 6 as the kernels
$S=\hat I$, $g=\delta K$ of M20; two-type extension M28; burst law reused in
M23 and M31. Specialises (dropping immigration and random catastrophe sizes)
the birth--immigration--catastrophe model of Brockwell, Gani \& Resnick
(1982). Master table of closed forms: Appendix \texttt{dist:app:formulae}.

% ---------------------------------------------------------------
\subsection{M7 --- Discrete-generation catastrophe variant}

\paragraph{Class.} Discrete-time variant of M6 used as a cautionary remark.

\paragraph{State space.} Generation-indexed population counts.

\paragraph{Dynamics.} Per generation, each individual independently
initiates catastrophe with probability $\chi\in(0,1)$; write
$\widetilde X_i$ for the catastrophe-free population and
$\mathcal E_n=\sum_{i=0}^{n-1}\widetilde X_i$ for cumulative exposure. Then
\[
\PR{\text{no catastrophe before generation }n}
=\EX{(1-\chi)^{\mathcal E_n}}
\qquad\text{(\texttt{m:eq:norupture})}.
\]

\paragraph{Principal results.} Jensen's inequality
$\EX{(1-\chi)^{\mathcal E_n}}\ge(1-\chi)^{\EX{\mathcal E_n}}$: a mean-field
hazard substitution \emph{underestimates} the non-catastrophe probability
(Remark \texttt{m:rem:notmeanfield}).

\paragraph{Location and connections.} \S2.2.2, Remark
\texttt{m:rem:notmeanfield}. Discrete-time shadow of M6's $\delta n$
hazard; motivates keeping the full distribution rather than substituting
means.

% ---------------------------------------------------------------
\subsection{M8 --- Time-inhomogeneous linear birth--death process}

\paragraph{Class.} CTMC with time-dependent rates (instance of F3).

\paragraph{State space.} $\{0,1,2,\ldots\}$; $0$ absorbing.

\paragraph{Dynamics.} Birth and death rates $\lambda(t),\mu(t)$ (per
capita, linear in the count, moving in time).

\paragraph{Parameters.} Arbitrary measurable rate schedules.

\paragraph{Initial conditions.} $X_0=n_0$.

\paragraph{Principal results.} Mean (\texttt{m:eq:meanodeinhom}):
\[
\EX{X_t\mid X_0=n_0}=n_0\exp\Bigl\{\int_0^t(\lambda(r)-\mu(r))\dd r\Bigr\};
\]
backward extinction Riccati for $u(s,T)=\PR{X_T=0\mid X_s=1}$
(\texttt{m:eq:riccatiext}):
\[
-\partial_su=\mu(s)-(\lambda(s)+\mu(s))u+\lambda(s)u^2,\qquad u(T,T)=0,
\]
with Kendall's forward-endpoint formula (\texttt{m:eq:inhomextforward})
\[
u_{\rm f}(t)=\frac{B(t)}{1+B(t)},\qquad
B(t)=\int_0^t\mu(s)\exp\Bigl\{\int_0^s(\mu(r)-\lambda(r))\dd r\Bigr\}\dd s .
\]

\paragraph{Location and connections.} \S2.2.3 (\texttt{m:sec:tibd}). Base
for the Chapter 9 programme (M10 and the extinction variant of M10); the
conclusion notes that ``what costs the generating function is not feedback
as such but continuous bending of a rate'' (Chapter 10). A named special
case (``subductive proliferation'', $\lambda(t)=\beta\ee^{-\alpha t}+\gamma$,
constant death) is mentioned in Chapter 8's discussion but not developed.

% ---------------------------------------------------------------
\subsection{M9 --- Seasonal Lotka--Volterra process}

\paragraph{Class.} Two-population CTMC with externally imposed periodic
rates (F3 instance).

\paragraph{State space.} $(x,y)$ = (prey, predator) counts,
$\mathbb Z_{\ge0}^2$.

\paragraph{Dynamics.} (\texttt{m:eq:lvrate}, \texttt{m:eq:lvkernel}):
\[
q_{(x,y),(x+1,y)}(t)=b\,\mathcal S^+(t)\,x,\qquad
q_{(x,y),(x-1,y+1)}(t)=e\,xy,\qquad
q_{(x,y),(x,y-1)}(t)=d\,y,
\]
\[
\mathcal S^+(t)=\max\{0,1+A\sin(\omega t)\},
\]
with a closed season of length $(2/\omega)\arccos(1/A)$ per period when
$A>1$.

\paragraph{Parameters.} $b$ prey birth, $e$ predation, $d$ predator death,
$A$ forcing amplitude, $\omega$ frequency (figure parameters
$A=2$, $\omega=2\pi$, $b=2$, $e=0.005$, $d=2$, $(300,100)$).

\paragraph{Principal results.} Exact marginal on the predator-free axis
(time-inhomogeneous Yule); mean-field ODE $\dot x=b\mathcal S^+(t)x-exy$,
$\dot y=exy-dy$ does not close ($\EX{XY}$ enters); simulation by thinning
against $\mathcal S^+(t)\le1+A$.

\paragraph{Location and connections.} \S2.2.3 (\texttt{m:sec:seasonallv}).
Worked example of externally forced rates; contrast class for the coupled
systems F4/M11/M12.

% ---------------------------------------------------------------
\subsection{M10 --- Logistic speciation model (inhomogeneous Yule)}
\label{sec:M10}

\paragraph{Class.} Non-homogeneous pure-birth CTMC driven by a
deterministic logistic environment (F3, coupling type (ii) of F4).

\paragraph{State space.} $S_t\in\{1,2,\ldots\}$ species count;
deterministic environment $N(t)\in(N_0,K)$.

\paragraph{Dynamics.} Environment ODE and rate law
(\texttt{m:eq:logisticode}, \texttt{var:eq:logistic}, \texttt{var:eq:rate1}):
\[
\dd N/\dd t=\varrho N\Bigl(1-\frac NK\Bigr),\quad N(0)=N_0,\quad
N(t)=\frac{K\ee^{\varrho t}}{\xi+\ee^{\varrho t}},\quad\xi=\frac K{N_0}-1,
\]
\[
\PR{S_{t+\Delta t}=S_t+1\mid S_t=n}=\beta(t)\,n\,\Delta t+o(\Delta t),
\qquad
\beta(t)=\sigma\bigl(K-N(t)\bigr)=\sigma K\frac{\xi}{\xi+\ee^{\varrho t}} .
\]
The coupling runs one way: the environment sets the speciation clock.

\paragraph{Parameters.} $\varrho$ growth rate, $K$ capacity, $N_0$ founding
abundance, $\sigma$ speciation coefficient; cumulative intensity
$\mathcal B(t)=\int_0^t\beta(u)\dd u=
\sigma K\{t-\varrho^{-1}\ln(\xi+\ee^{\varrho t})+\varrho^{-1}\ln(\xi+1)\}$.

\paragraph{Initial conditions.} $S_0=1$.

\paragraph{Principal results.} Geometric marginal at every time --- the
Yule law with $\lambda t$ replaced by $\mathcal B(t)$ (Proposition
\texttt{var:prop:yule-geom}):
\[
\PR{S_t=n}=\ee^{-\mathcal B(t)}\bigl(1-\ee^{-\mathcal B(t)}\bigr)^{n-1},
\qquad \EX{S_t}=\ee^{\mathcal B(t)};
\]
limit law under the logistic rate (Corollary 9.3
= \texttt{var:cor:logistic-limit}):
\[
\mathcal B(\infty)=\frac{\sigma K}{\varrho}\ln\frac K{N_0},\qquad
\EX{S_\infty}=\Bigl(\frac K{N_0}\Bigr)^{\sigma K/\varrho},
\]
with only the dimensionless groups $K/N_0$ and $\sigma K/\varrho$
surviving; pgf $G(z,t)=z/\{z+(1-z)\ee^{\mathcal B(t)}\}$ via the method of
characteristics (Appendix \texttt{var:app:pgf}).

\paragraph{Variants.} Alternative rate laws (\texttt{var:eq:altrates}):
$\beta_2=\sigma(K-N)+\eta$ and $\beta_4=\sigma\,\dd N/\dd t+\eta$ (no limit
law, $\mathcal B\sim\eta t$); $\beta_3=\sigma\,\dd N/\dd t$ with
$\mathcal B(\infty)=\sigma(K-N_0)$; $\beta_5=\sigma(K_t-N(t))$ with randomly
switching capacity --- mixture of geometrics, open. Extinction variant
(\texttt{var:sec:withextinction}): add constant per-species death $\mu$;
$\EX{S_t}=\exp\{\mathcal B(t)-\mu t\}$ turns over where $\beta(t)=\mu$ and
tends to $0$; no limit law.

\paragraph{Location and connections.} Mean-only version in \S2.2.3
(\texttt{m:sec:logisticspeciation}); full distribution theory Chapter 9
(\texttt{var:sec:speciation}). Constant-rate limit is M4; positions the
``push of the past'' debate (recovered only through a moving rate).

% ---------------------------------------------------------------
\subsection{M11 --- Accumulation with gated release}

\paragraph{Class.} PDMP (F4, coupling type (iii): two-way).

\paragraph{State space.} Continuous $(A_t,R_t)\in\mathbb R_{\ge0}^2$
(accumulated and transferred amounts) and gate $G_t\in\{0,1\}$.

\paragraph{Dynamics.} Flow between jumps and jump rates
(\texttt{m:eq:gatedflow}, \texttt{m:eq:gatedjumps}):
\[
\dd A_t/\dd t=\rho-G_tcA_t,\qquad \dd R_t/\dd t=G_tcA_t,
\]
\[
\PR{G_{t+\Delta t}=1\mid G_t=0,A_t=a}=\eta a\,\Delta t+o(\Delta t),
\qquad
\PR{G_{t+\Delta t}=0\mid G_t=1}=\gamma\,\Delta t+o(\Delta t).
\]

\paragraph{Parameters.} $\rho$ inflow, $c$ transfer coefficient, $\eta$
(per amount) opening rate, $\gamma$ closing rate; figure
$(\rho,c,\eta,\gamma)=(2,4,0.015,1.45)$, $(A_0,R_0,G_0)=(5,0,0)$.

\paragraph{Principal results.} Pathwise balance $A_t+R_t=A_0+R_0+\rho t$;
first moment unclosed ($\frac{\dd}{\dd t}\EX{A_t}=\rho-c\EX{G_tA_t}$);
simulation solves a quadratic cumulative hazard
$\eta(A_tr+\tfrac12\rho r^2)$ against an $\Exp(1)$ variate.

\paragraph{Location and connections.} \S2.2.4 (\texttt{m:sec:gated}).
Single-cell caricature of intermittent (quasi-bursting) release, cf.\
Chapter 6's kernels and M25.

% ---------------------------------------------------------------
\subsection{M12 --- Accumulation with surging gated release}

\paragraph{Class.} PDMP (F4, coupling type (i): autonomous mode chain).

\paragraph{State space.} As M11.

\paragraph{Dynamics.} Same flow; autonomous gate with
$\PR{G_{t+\Delta t}=1\mid G_t=0}=\nu\Delta t+o(\Delta t)$
(\texttt{m:eq:surgingjumps}); $G_t$ is an autonomous two-state CTMC.

\paragraph{Parameters.} $\nu$ opening, $\gamma$ closing rates; others as
M11.

\paragraph{Principal results.} $\frac{\dd}{\dd t}\EX{G_t}=\nu(1-\EX{G_t})-
\gamma\EX{G_t}$, so $\EX{G_t}\to\nu/(\nu+\gamma)$; balance law of M11
unchanged; population moments still contain $\EX{G_tA_t}$.

\paragraph{Location and connections.} \S2.2.4 (\texttt{m:sec:surging}).
Contrast class to M11: autonomy restores some closure but not the joint
moment.

% ---------------------------------------------------------------
\subsection{M13 --- Absorption-only process}

\paragraph{Class.} Bivariate CTMC (immigration--death structure).

\paragraph{State space.} $(X_t,Y_t)$ = (inside, outside) a compartment,
$\mathbb Z_{\ge0}^2$; nothing returns from the compartment.

\paragraph{Dynamics.} (\texttt{m:app:absonly}):
$\PR{\Delta X=1,\Delta Y=-1\mid Y_t=j}=j\alpha\Delta t$ (absorption), no
other events.

\paragraph{Parameters.} $\alpha$ per-particle uptake rate.

\paragraph{Initial conditions.} $X_0=0$, $Y_0=N$.

\paragraph{Principal results.} $X_t\sim\operatorname{Binomial}(N,
1-\ee^{-\alpha t})$; $G(x,y,t)=\{(1-\ee^{-\alpha t})x+\ee^{-\alpha t}y\}^N$
(\texttt{m:eq:Gabsonly}).

\paragraph{Location and connections.} Appendix 2.B
(\texttt{m:app:absonly}). Simplest member of the absorption family
M13--M15, which model transfer from medium into a compartment after a
rupture.

% ---------------------------------------------------------------
\subsection{M14 --- Absorption--death process}

\paragraph{Class.} Bivariate CTMC; the worked example of the method of
characteristics.

\paragraph{State space.} As M13.

\paragraph{Dynamics.} (\texttt{m:sec:mocworked}):
\[
\PR{\Delta X=1,\Delta Y=-1\mid Y_t=j}=j\alpha\Delta t,\qquad
\PR{\Delta X=-1,\Delta Y=0\mid X_t=i}=i\mu\Delta t .
\]

\paragraph{Parameters.} $\alpha$ uptake, $\mu$ interior death.

\paragraph{Initial conditions.} $X_0=0$, $Y_0=N$.

\paragraph{Principal results.} Forward equations (\texttt{m:eq:FK2});
pgf PDE $\partial_tG=\mu(1-x)\partial_xG+\alpha(x-y)\partial_yG$,
$G(x,y,0)=y^N$ (\texttt{m:eq:PDE2}); closed-form solution
(\texttt{m:eq:Gabsdeath}):
\[
G(x,y,t)=\Bigl(1-\tfrac{\mu}{\mu-\alpha}\ee^{-\alpha t}
+\tfrac{\alpha}{\mu-\alpha}\ee^{-\mu t}
+\tfrac{\alpha}{\mu-\alpha}(\ee^{-\alpha t}-\ee^{-\mu t})x
+\ee^{-\alpha t}y\Bigr)^{N},
\]
recovered exactly by the characteristic system
(\texttt{m:eq:ch11}--\texttt{m:eq:ch14}).

\paragraph{Location and connections.} \S2.3.3 (\texttt{m:sec:mocworked}).
Motivates Appendix 2.B: adding interior birth makes the $x$-characteristic
Riccati, i.e.\ M15.

% ---------------------------------------------------------------
\subsection{M15 --- Absorption--birth--death process [thesis result]}

\paragraph{Class.} Bivariate CTMC; the hardest member of the absorption
family; closed form is a result of this thesis (Proposition
\texttt{m:prop:Gabd}).

\paragraph{State space.} As M13.

\paragraph{Dynamics.} (\texttt{m:app:absbirthdeath}): uptake
$\Delta X=1,\Delta Y=-1$ at rate $j\alpha$; interior birth $\Delta X=1$ at
rate $i\lambda$; interior death $\Delta X=-1$ at rate $i\mu$.

\paragraph{Parameters.} $\alpha$ uptake; $\lambda,\mu$ interior BD rates;
$b_1=\lambda-\mu$ (net interior growth); $\omega=\alpha/b_1$ (absorption in
units of net growth).

\paragraph{Initial conditions.} $X_0=0$, $Y_0=N$.

\paragraph{Principal results.} PDE (\texttt{m:eq:PDE3})
\[
\partial_tG=\bigl(\mu-(\mu+\lambda)x+\lambda x^2\bigr)\partial_xG
+\alpha(x-y)\partial_yG ,
\]
whose $x$-coefficient is exactly the birth--death quadratic of M5 with
roots $x_+=1$, $x_-=\mu/\lambda$. With
$W=(x-x_+)/(x-x_-)$ and the series function
\[
\Psi(\zeta)=x_++\frac{\alpha}{\lambda}\sum_{n=1}^{\infty}\frac{\zeta^n}{n+\omega},
\qquad \omega=\frac{\alpha}{\lambda-\mu},
\]
the boxed solution (Proposition \texttt{m:prop:Gabd}) is
\[
G(x,y,t)=\Bigl\{\ee^{-\alpha t}\bigl[y-\Psi(W)\bigr]
+\Psi\bigl(W\ee^{b_1t}\bigr)\Bigr\}^{N}.
\]
Hypergeometric form $\Psi(\zeta)=x_++(x_+-x_-)\{{}_2F_1(1,\omega;\omega+1;
\zeta)-1\}$; parameter-regular representation via a convolution with the
interior survival pgf (finite at $\lambda=\mu$, where the boxed form
degenerates); resonance poles at $\alpha=n(\mu-\lambda)$ cancel to a secular
factor $t\ee^{-\alpha t}$; coefficient recovery $p_{i,j}(t)
=\binom Nj\ee^{-\alpha tj}[x^i]h(x,t)^{N-j}$. Verified numerically to
$\sim10^{-13}$.

\paragraph{Location and connections.} Appendix 2.B
(\texttt{m:app:absbirthdeath}), supported by the integral identity of
Appendix 2.C (\texttt{m:prop:hypIden}) and the extraction machinery of
Appendix 2.D. Connects M5's quadratic, M14's characteristics, and the
rupture-transfer question of the biology chapters.

% ============================================================================
\section{Catalogue, part II --- distribution theory and pathogenesis (Chapters 5--8)}
% ============================================================================

% ---------------------------------------------------------------
\subsection{M16 --- Chained immediate-transfer process ($\mu=0$)}
\label{sec:M16}

\paragraph{Class.} Chained sequence of single-cell birth--catastrophe
CTMCs (M6 at $\mu=0$) coupled by wholesale transfer of the burst into the
next fresh cell; ``no extracellular phase, no clearance, no re-infection''.

\paragraph{State space.} Per stage: load $n\in\{1,2,\ldots\}$; across
stages: rupture sizes $r_k$, rupture times $t_k$, inter-rupture intervals
$T_k$.

\paragraph{Dynamics.} From load $n$, next event is birth w.p.\
$\lambda/(\lambda+\delta)=\rho$ and catastrophe w.p.\
$s=\delta/(\lambda+\delta)$, independent of $n$ (the load ``drops out'');
holding time $\Exp((\lambda+\delta)n)$. Births before catastrophe
$G\sim\Geom_0(s)$, $\PR{G=g}=s\rho^g$, independent across cells; rupture
size decomposition (\texttt{dist:eq:rkdecomp})
\[
r_k=1+G_1+\cdots+G_k .
\]

\paragraph{Parameters.} $\lambda,\delta$ ($\mu=0$); $s,\rho$ as above.

\paragraph{Initial conditions.} One cell, one particle, $t_0=0$.

\paragraph{Principal results.} Negative-binomial rupture law
(\texttt{dist:eq:rklaw}):
\[
\PR{r_k=n}=\binom{n+k-2}{k-1}s^k\rho^{\,n-1},\qquad n\ge1,
\]
with $\EX{r_k}=1+k\lambda/\delta$, $\Var{r_k}=k\rho/s^2$, pgf
$\EX{z^{r_k}}=z\{s/(1-\rho z)\}^k$; first rupture from $m$ founders:
$\PR{t_1>t}=I(t)^m$; inter-rupture Laplace transform in closed
hypergeometric form (\texttt{dist:eq:ltkhyp}) with mean
$\EX{T(k)}=\sum_{m\ge0}\rho^m/\{(\lambda+\delta)(k+m)\}$, strictly
decreasing in $k$. Exact simulation by coin flips and exponential waits;
verified 28/28 (\texttt{dist:app:verification}). The structure breaks for
$\mu>0$ (a cell then has two ways to stop); joint laws of $(t_k,r_k)$ at
$\mu>0$ remain open.

\paragraph{Location and connections.} Chapter 5 \S10
(\texttt{dist:sec:chained}); reused in Chapter 8
(\texttt{path:sec:chainedruptures}) as the ``release schedule alone''
model. Interpolates between single-cell theory (M6) and populations (M19).

% ---------------------------------------------------------------
\subsection{M17 --- Budding comparator cell}

\paragraph{Class.} Two independent Poisson streams on one cell (CTMC on
released counts crossed with cell state).

\paragraph{State space.} Cumulative releases $k\in\{0,1,2,\ldots\}$ while
the cell is alive.

\paragraph{Dynamics.} Release of one particle at rate $p$; cell death at
rate $d_I$; the streams run until death.

\paragraph{Parameters.} $p$ release rate; $d_I$ infected-cell death rate
(mean lifetime $1/d_I$).

\paragraph{Initial conditions.} Cell infected at $t=0$.

\paragraph{Principal results.} (\texttt{dist:eq:bud}, \texttt{p:eq:bud-cell}):
\[
I_{\rm bud}(t)=\ee^{-d_It},\qquad
V_{\rm bud}(t)=\frac{p}{d_I}\bigl(1-\ee^{-d_It}\bigr),\qquad
\PR{K_{\rm bud}=k}=\frac{d_I}{d_I+p}\Bigl(\frac{p}{d_I+p}\Bigr)^{k},
\]
geometric on $\{0,1,2,\ldots\}$ (releases before the terminating death).

\paragraph{Location and connections.} Chapter 5 \S11
(\texttt{dist:sec:budding}); Chapter 6 \S\texttt{p:sec:budding}, where it
is shown that the BMVR (M18) describes a budding pathogen \emph{exactly};
matched-mean comparator in M23 and M31. Structural opposite of M6: release
and death independent vs.\ identical events.

% ---------------------------------------------------------------
\subsection{M18 --- Basic model of viral replication (BMVR) \quad [D]}

\paragraph{Class.} \textbf{[D]} Deterministic ODE comparator (not Markov).

\paragraph{State variables.} Target cells $T$, infected cells $I$, free
virions $V$.

\paragraph{Equations.} Three-variable form (\S1.6):
\[
\dd T/\dd t=s-\gamma TV,\qquad
\dd I/\dd t=\gamma TV-d_II,\qquad
\dd V/\dd t=pI-cV,
\]
and the constant-target reduction carried through Chapter 6
(\texttt{p:eq:bmvr-classical}):
\[
\dd I/\dd t=\gamma TV-d_II,\qquad \dd V/\dd t=pI-cV .
\]

\paragraph{Parameters.} $\gamma$ infection rate, $T$ fixed target density,
$d_I$ infected-cell death, $p$ release rate, $c$ clearance; basic
reproduction number $R_0=\gamma Tp/(c\,d_I)$ (\texttt{p:eq:R0classical}).

\paragraph{Role.} The model Chapter 6 replaces: constants $p,d_I$ are
exchanged for the age-dependent kernels of M20, and the BMVR is recovered
exactly as the exponential-phase projection (M21). Assumptions challenged
by the thesis: one-by-one release at constant rate; release and death
separate independent events (budding).

\paragraph{Location.} Chapter 1 \S1.6; Chapter 6
\S\texttt{p:sec:classical}.

% ---------------------------------------------------------------
\subsection{M19 --- Stochastic within-host population process}

\paragraph{Class.} Branching-type CTMC population model (linear rates);
the ``first-moment-exactness'' model.

\paragraph{State space.} Product of per-cell ages/loads (each an M6
process) and free-particle count $V_t$.

\paragraph{Dynamics.} Each infection event starts an independent BDC cell
(M6); each released particle independently infects at rate $\gamma T$ or
clears at rate $c$; the $I_0$ founding cells have age zero. Infection
events form a point process with compensator $\gamma T\,V^{\omega}(t)\dd t$
(Campbell's formula). Gillespie specification (\texttt{p:fig:gillespie}):
``birth--death--catastrophe cells with impulse releases at burst times,
exponential clearance, no reinfection''.

\paragraph{Parameters.} $(\lambda,\mu,\delta)$ per cell; $\gamma T$, $c$;
$I_0$.

\paragraph{Initial conditions.} $I_0$ cells of age zero (assumption A3);
$V_0$ implicit.

\paragraph{Principal results.} Proposition \texttt{p:prop:first-moment}:
its expectations satisfy the renewal system M20 \emph{exactly} --- linear
incidence (A1) is what closes the moment hierarchy; ensemble Gillespie
means match the renewal means to 1--2\% (62-check verification suite).

\paragraph{Location and connections.} Chapter 6
\S\texttt{p:sec:renewal}. Stochastic ground truth under M20; nonlinear
incidence (M27) breaks the exactness.

% ---------------------------------------------------------------
\subsection{M20 --- Burst-aware renewal system and the $(S,g)$ skeleton
\quad [D: exact mean equations of M19]}

\paragraph{Class.} Deterministic renewal/Volterra system
(McKendrick--von Foerster age structure); exact first moments of M19.

\paragraph{Equations.} Definition \texttt{p:def:renewal-system}: with
incidence $i(t)=\gamma TV(t)$ and the BDC kernels
$S=\hat I$ (survival) and $g=\delta K$ (release rate),
\[
I(t)=I_0\,\hat I(t)+(i*\hat I)(t),
\qquad
\dd V/\dd t=I_0\,g(t)+(i*g)(t)-cV(t)
\qquad\text{(\texttt{p:eq:Icell}, \texttt{p:eq:Vfree})}.
\]
General skeleton (\texttt{p:eq:skeleton}), any pair $(S,g)$:
$I=I_0S+(i*S)$, $V'=I_0g+(i*g)-cV$, with
$R_0=(\gamma T/c)\int_0^\infty g$ under linear incidence.

\paragraph{Instances of the skeleton} (Table \texttt{p:tab:sg}): standard
BMVR ($S=\ee^{-d_I\alpha}$, $g=p\ee^{-d_I\alpha}$); eclipse model;
HCV multiscale (Guedj / Rong--Perelson); absorbing BDC (this chapter);
MOI-$k$ BDC ($I^k-D^k$, $\delta K_k$); reset catastrophe (M25); two-type /
GATE (Chapter 7). Exponential kernels reduce it to M18 (checked to
$5\times10^{-5}$).

\paragraph{Location and connections.} Chapter 6
\S\texttt{p:sec:renewal}, \S\texttt{p:sec:generality}; kernels imported
from M6/M15-level theory (Chapters 4--5).

% ---------------------------------------------------------------
\subsection{M21 --- Effective-parameter projection \quad [D: exact map]}

\paragraph{Class.} \textbf{[D]} Exact map from M20 to constant-parameter
M18 along exponential phases $i(t)=i_0\ee^{rt}$
(Theorem \texttt{p:thm:projection}).

\paragraph{Definition.} (\texttt{p:eq:peff}):
\[
p_{\rm eff}(r)=\frac{\delta\,\Lap K(r)}{\Lap{\hat I}(r)},
\qquad
d_{I,\rm eff}(r)=\frac{1}{\Lap{\hat I}(r)}-r,
\qquad r>-\theta,
\]
returning the characteristic equation $r+c=\gamma T\,\delta\Lap K(r)$
(\texttt{p:eq:char}).

\paragraph{Ends and properties.} $r=0$: $p_{\rm eff}(0)=V_\infty/\EX{T_{\rm
prod}}$, $d_{I,\rm eff}(0)=\lambda/\log\frac{a}{a-1}$; $r\to\infty$:
$p_{\rm eff}\to\delta$, $d_{I,\rm eff}\to\mu+\delta$; old-cell limit
$p_{\rm eff}(-\theta)\to\delta\EX{X^2}_{\rm QS}$ (Proposition
\texttt{p:prop:oldcell}); span $p_{\rm eff}(-\theta)/p_{\rm eff}(\infty)
=\EX{X^2}_{\rm QS}$ (231 at $(1,0,0.1)$); monotonicity via
$\frac{\dd}{\dd r}\log p_{\rm eff}=\EX{\alpha}_{\pi_{S,r}}-\EX{\alpha}_
{\pi_{g,r}}$ (numerically supported at 4000 triples). Identifiability
fails: $(\lambda,\mu,\delta)\mapsto(p_{\rm eff},d_{I,\rm eff})$ is not
invertible; candidate third observables are the geometric ratio $1/a$, the
clearance fraction $b$, and the shape of $\delta J$.

\paragraph{Location.} Chapter 6 \S\texttt{p:sec:projection},
\S\texttt{p:sec:identifiability}; hypergeometric Laplace transforms in
Appendix \texttt{p:app:transforms}.

% ---------------------------------------------------------------
\subsection{M22 --- Generation kernel and $R_0$ \quad [D: Euler--Lotka functional]}

\paragraph{Class.} \textbf{[D]} Functional of the release kernel $g$.

\paragraph{Definition.} (\texttt{p:eq:genkernel}):
\[
\mathcal A(\alpha)=\gamma T\int_0^\alpha g(u)\ee^{-c(\alpha-u)}\dd u,
\qquad \Lap{\mathcal A}(r)=\frac{\gamma T\,\Lap g(r)}{r+c},
\]
$\Lap{\mathcal A}(r)=1$ the characteristic equation; $\mathcal A/R_0$ the
generation-time density.

\paragraph{Principal results.} (\texttt{p:eq:R0gen})
\[
R_0=\int_0^\infty\mathcal A(\alpha)\dd\alpha
=\frac{\gamma T\,V_\infty}{c},
\]
and Theorem \texttt{p:thm:R0inv}: $R_0$ is unchanged by bursting at fixed
total output per cell --- the schedule becomes visible only through
establishment (M23) or a saturating defence (M31).

\paragraph{Location.} Chapter 6 \S\texttt{p:sec:generation-kernel};
continuum counterpart of the discrete branching mean $m=qV_\infty$ in M23.

% ---------------------------------------------------------------
\subsection{M23 --- Early-infection establishment branching process (bursting vs budding)}

\paragraph{Class.} Galton--Watson branching process of infected cells;
offspring distributions in closed form for both mechanisms.

\paragraph{Dynamics.} Each released particle independently infects with
probability
\[
q=\frac{\gamma T}{\gamma T+c}
\qquad\text{(\texttt{p:eq:q})}
\]
or is cleared. Bursting offspring pgf: thinning the M6 burst pgf
$G_K(z)=b+\frac{\delta}{\lambda}\frac z{a-z}$ at $y=1-q+qz$
(\texttt{p:eq:Goff}):
\[
G_{\rm off}(z)=b+\frac{\delta}{\lambda}\,\frac{y}{a-y},\qquad y=1-q+qz ,
\]
with mean $m=qV_\infty$. Budding comparator: geometric offspring on
$\{0,1,2,\ldots\}$ with the same mean (matched at $p/d_I=V_\infty$).

\paragraph{Parameters.} $q$; the M6-derived $a,b,V_\infty$; flooding
parameter (Definition \texttt{p:def:L})
\[
L:=a(1-b)=a-\frac{\mu}{\lambda},\qquad V_\infty=\frac{L}{a-1}.
\]

\paragraph{Principal results.} Extinction probabilities for $m>1$
(\texttt{p:eq:zext}):
\[
z_{\rm ext}^{\rm burst}=\frac{a(1-q+qb)-1+q}{q},\qquad
z_{\rm ext}^{\rm bud}=\frac1m,
\]
and the flooding theorem (Theorem \texttt{p:thm:flood},
\texttt{p:eq:flood}):
\[
z_{\rm ext}^{\rm burst}-z_{\rm ext}^{\rm bud}=(L-1)\Bigl(\frac1m-1\Bigr),
\]
so bursting establishes more easily than budding iff $L>1$ iff
$\delta(\lambda+\mu)>\lambda\mu$ (\texttt{p:eq:floodcrit}; convention: holds
for all $\delta>0$ when $\mu=0$). At $L=1$ the two offspring laws coincide
exactly (geometric; Proposition \texttt{p:prop:L1}), and
$\Var_{\rm bud}-\Var_{\rm burst}=2q^2V_\infty(L-1)/(a-1)$
(\texttt{p:eq:varorder}).

\paragraph{Location and connections.} Chapter 6
\S\texttt{p:sec:flooding}. Built on M6's burst law; the boundary curve
$\delta_*=\lambda\mu/(\lambda+\mu)$ in the rate plane; the companion
invasion-speed comparison (M24): $r_{\rm bud}>r_{\rm burst}$ at matched
$R_0=2$ via the Laplace-transform order $\Lap g_{\rm burst}(r)\le
p/(r+d_I)$ for all $r>0$ --- a Pareto trade-off with the extinction
probability.

% ---------------------------------------------------------------
\subsection{M24 --- Invasion-speed comparison \quad [D]}

\paragraph{Class.} \textbf{[D]} Deterministic comparison of Malthusian
growth rates at matched $R_0$; not a process.

\paragraph{Setting.} Budding: eigenvalue relation
$(r+d_I)(r+c)=\gamma Tp$ (from M18); bursting: the characteristic equation
of M20/M21, $r+c=\gamma T\,\delta\Lap K(r)$; matched at $R_0=2$,
$d_I=1/\EX{T_{\rm prod}}$, $p=V_\infty d_I$ ($=p_{\rm eff}(0)$).

\paragraph{Principal results.} Ordering (\texttt{p:eq:rorder}):
$r_{\rm bud}>r_{\rm burst}$ in all three named BDC regimes (e.g.\ $0.250$
vs $0.180$ at $(1,0,0.1)$); Proposition \texttt{p:prop:rorder} reduces it
to the Laplace-transform order
\[
\Lap g_{\rm burst}(r)\ \le\ \Lap g_{\rm bud}(r)=\frac{p}{r+d_I}
\qquad\text{for all } r>0
\]
(checked at 1500 triples; the usual stochastic order does \emph{not}
hold). Mechanism: the generation-time density $\mathcal A/R_0$ sits earlier
for budding; varying $\delta$ traces a Pareto frontier against the
extinction probability of M23 (\texttt{p:fig:pareto}).

\paragraph{Location and connections.} Chapter 6
\S\texttt{p:sec:tradeoff}. Deterministic counterpart of M23: budding wins
on speed, bursting wins on establishment when $L>1$.

% ---------------------------------------------------------------
\subsection{M25 --- Reset process (recurrent partial release)}

\paragraph{Class.} Single-cell CTMC with non-absorbing dumps and a separate
cell-death clock.

\paragraph{State space.} Load $n\in\{0,1,2,\ldots\}$ (may re-accumulate)
$\times$ cell alive/dead.

\paragraph{Dynamics.} While alive: birth $\lambda n$, death $\mu n$,
optional immigration $\nu$, and a dump at rate $\delta n$ that releases the
whole pool and returns the load to zero; cell dies independently at rate
$d_I$ (\S\texttt{p:sec:reset}).

\paragraph{Parameters.} $\lambda,\mu,\delta,(\nu),d_I$.

\paragraph{Principal results.} Kernels (\texttt{p:eq:reset-kernels}):
\[
S(\alpha)=\ee^{-d_I\alpha},\qquad
g(\alpha)=\delta\,\ee^{-d_I\alpha}\,\EX{X_\alpha^2},
\]
effective parameters $p_{\rm eff}(r)=(r+d_I)\Lap g(r)$, $d_{I,\rm eff}
\equiv d_I$; $R_0$ finite only for $d_I>0$; mature-cell limit
$p_\infty=\delta\,\EX_\pi[X^2]$ if a stationary law $\pi$ exists
(\texttt{p:eq:reset-mature}). Intermediate between M17 and M6 on the M20
skeleton.

\paragraph{Location and connections.} Chapter 6 \S\texttt{p:sec:reset};
``multi-pulse caricature between full dumping and continuous export''.

% ---------------------------------------------------------------
\subsection{M26 --- Partial-release stock model \quad [D: mean-field]}

\paragraph{Class.} \textbf{[D]} Deterministic mean-field ODE system;
boolean thinning of releases (each unit released independently w.p.\
$\varphi$), event clock $\delta X$.

\paragraph{Equations.} (\texttt{p:eq:partial}), with $Q$ = total
intracellular stock across productive cells:
\[
\dd I/\dd t=\gamma TV-d_II,\qquad
\dd Q/\dd t=\nu I-\delta\varphi\,\frac{Q^2}{I}-d_IQ,\qquad
\dd V/\dd t=\delta\varphi\,\frac{Q^2}{I}-cV .
\]

\paragraph{Principal observation.} With a load-independent clock
($\lambda(X)=\delta$) the mean export is $\delta\varphi Q$ and partial
release is invisible in the mean-field equations; it becomes visible when
the clock is load-dependent, and is always visible in single-cell release
distributions. Used to map antiviral modes of action
($\downarrow\lambda$, $\uparrow\mu$, $\downarrow\gamma T$, $\downarrow
\varphi=1-\epsilon_s$).

\paragraph{Location.} Chapter 6 \S\texttt{p:sec:partial}.

% ---------------------------------------------------------------
\subsection{M27 --- HIV skeleton variant \quad [specification]}

\paragraph{Class.} Specification on the M20 skeleton; not solved.

\paragraph{Dynamics.} Kernels $S_{\rm HIV}$ (productive-lineage survival)
and $g_{\rm HIV}$ (productive-sojourn density) with a multi-stage (Erlang,
$n\approx5$) eclipse delay; intracellular load approximately
immigration-like ($\bar x'\approx\nu-\mu\bar x$); incidence with Allee
effect (\texttt{p:eq:allee}):
\[
f(V,T)=\gamma TV\cdot\frac{V}{K_A+V}.
\]

\paragraph{Consequences noted.} $R_0$ remains the large-population
threshold; establishment additionally governed by $K_A$; the
first-moment exactness of M19/M20 no longer holds (nonlinear incidence
breaks moment closure); the independent-branching approximation fails near
HIV's observed threshold ($\sim$2.3 founder cells, Hataye et al.).

\paragraph{Location.} Chapter 6 \S\texttt{p:sec:hiv}.

% ---------------------------------------------------------------
\subsection{M28 --- Two-type birth--death--conversion process with global catastrophe}
\label{sec:M28}

\paragraph{Class.} Two-type continuous-time branching process (triangular:
type 2 autonomous) with one-way conversion and a single globally absorbing
killing event.

\paragraph{State space.} $(x,y)\in\mathbb Z_{\ge0}^2$ (type-1 and type-2
counts) plus the absorbing state $(\Rst,\Rst)$.

\paragraph{Dynamics.} Transition table (\S\texttt{sec:model}):
\[
\begin{array}{rcll}
(x,y)\to(x+1,y) & \text{at rate} & \lambda_1x & \text{(type-1 birth)}\\
(x,y)\to(x-1,y) & \text{at rate} & \mu_1x & \text{(type-1 death)}\\
(x,y)\to(x-1,y+1) & \text{at rate} & \nu x & \text{(conversion, one-way)}\\
(x,y)\to(x,y+1) & \text{at rate} & \lambda_2y & \text{(type-2 birth)}\\
(x,y)\to(x,y-1) & \text{at rate} & \mu_2y & \text{(type-2 death)}\\
(x,y)\to(\Rst,\Rst) & \text{at rate} & \delta_1x+\delta_2y
& \text{(catastrophe, \texttt{eq:catastrophe-rate})}
\end{array}
\]
with $(\Rst,\Rst)$ absorbing.

\paragraph{Parameters.} $\lambda_i,\mu_i$ per-type birth/death; $\nu$
irreversible conversion; $\delta_1,\delta_2\ge0$ type-specific catastrophe
rates. Biological mapping (\texttt{tab:bio-mapping}): type 1 = early
intracellular phenotype, type 2 = adapted phenotype, conversion =
physiological adaptation, catastrophe = loss of macrophage containment.
Four labelled regimes (\S7 biology): EQ ($\delta_1=\delta_2$), MAT
($0<\delta_1<\delta_2$), GATE ($\delta_1=0<\delta_2$, $S'(0)=0$), EARLY
($\delta_1>\delta_2$).

\paragraph{Initial conditions.} Single type-dependent ancestors:
$S(t)=\PR_{(1,0)}\{\tau_c>t\}$, $G(t)=\PR_{(0,1)}\{\tau_c>t\}$,
$S(0)=G(0)=1$.

\paragraph{Principal results.}
\begin{itemize}
\item \emph{Feynman--Kac representation} (Proposition
  \texttt{prop:feynman-kac}): $S,G$ are occupation-time transforms of the
  catastrophe-free skeleton --- e.g.
  $S(t)=\EX_{(1,0)}\exp\{-\int_0^t(\delta_1X_s+\delta_2Y_s)\dd s\}$ ---
  the catastrophe clock being an inhomogeneous Poisson process of rate
  $\delta_1X_s+\delta_2Y_s$ along skeleton paths.
\item \emph{Triangular backward system} (\texttt{eq:backward-unscaled}),
  $S'(0)=-\delta_1$, $G'(0)=-\delta_2$:
  \[
  \dot S=\lambda_1S^2-(\lambda_1+\mu_1+\nu+\delta_1)S+\mu_1+\nu G,
  \qquad
  \dot G=\lambda_2G^2-(\lambda_2+\mu_2+\delta_2)G+\mu_2 .
  \]
\item \emph{Type-2 solution in closed form} (Proposition
  \texttt{prop:G-solution}): with roots
  $r_{2,\pm}=\frac{\hat\lambda_2+\hat\mu_2+\hat\delta_2\pm\hat\Delta_2}
  {2\hat\lambda_2}$, $\hat\Delta_2=\sqrt{(\hat\lambda_2+\hat\mu_2+
  \hat\delta_2)^2-4\hat\lambda_2\hat\mu_2}$ (hats = rates scaled by
  $\lambda_1$), $z_2(\hat t)=z_{2,0}\ee^{\hat q_2\hat t}$,
  \[
  \widehat G(\hat t)=r_{2,-}-\alpha_2\frac{z_2(\hat t)}{1-z_2(\hat t)}
  \qquad\Bigl(\alpha_2=\frac{\hat\Delta_2}{\hat\lambda_2},\ \
  \hat q_2=-\hat\Delta_2\Bigr).
  \]
\item \emph{Main theorem} (\texttt{tt:thm:main}; the chapter's headline
  result): after the shift $X=\widehat S-h$ by the lower root $h$ of
  $h^2-\hat\kappa_1h+\hat\mu_1+\hat\nu r_{2,-}=0$, the Riccati linearises
  to a Gauss hypergeometric equation with parameters
  \[
  C=1+\frac{\hat\kappa_1-2h}{\hat q_2},\qquad
  A+B=C-1=\frac{\hat\kappa_1-2h}{\hat q_2},\qquad
  AB=\frac{\hat\nu\alpha_2}{\hat q_2^2},
  \]
  and (basis $U(z)={}_2F_1(A,B;C;z)$, $V(z)=(-z)^{1-C}{}_2F_1(A-C+1,
  B-C+1;2-C;z)$ on the negative axis, mix fixed by
  $\Phi'(z_{2,0})/\Phi(z_{2,0})=L=(h-1)/(\hat q_2z_{2,0})$):
  \[
  \boxed{\ \widehat S(\hat t)=h-\hat q_2z_2(\hat t)\,
  \frac{K_UU'(z_2(\hat t))+K_VV'(z_2(\hat t))}
  {K_UU(z_2(\hat t))+K_VV(z_2(\hat t))}\ }
  \qquad\text{(\texttt{eq:S-final})}.
  \]
\item \emph{Long-time limits} (Corollary \texttt{cor:long-time}):
  $G\to r_{2,-}$, $S\to h$.
\item \emph{Boundary degenerations} (\texttt{app:limits}):
  $\delta_2=0$ gives an elementary Riccati solution with
  $R_0=(1-s_-)/(1-s_+)$ (Proposition \texttt{prop:delta2-zero});
  $\lambda_1=0$ linearises the type-1 equation; $\lambda_2=0$ reduces to
  Bessel functions $J_\rho,Y_\rho$ (Proposition
  \texttt{prop:lambda2-zero}); $\nu=0$ decouples into two one-type
  problems; equal rates restrict algebraically to the diagonal.
\item \emph{Killed (defective) bivariate pgf} (Appendix
  \texttt{app:catastrophe-pgf}, Proposition
  \texttt{prop:killed-pgf-backward}):
  $\mathcal F_i(x,y,t)=\EX_i[x^{X_t}y^{Y_t}\Ind\{\tau_c>t\}]$ satisfies
  the same backward system (with $\mathcal F_1=x$, $\mathcal F_2=y$ at
  $t=0$), solved by the same hypergeometric machinery with starting
  coordinate $z_{2,0}(y)=\frac{y-r_{2,-}}{y-r_{2,+}}$. Derived objects:
  absorbing-state marker $\mathcal H_i=\mathcal F_i+\rho[1-\mathcal F_i
  (1,1,t)]$, conditional pgf $\mathcal Q_i=\mathcal F_i/\mathcal F_i(1,1,t)$,
  flux GF $\mathcal J_i=\delta_1x\,\partial_x\mathcal F_i+\delta_2y\,
  \partial_y\mathcal F_i$, three conditional count laws, and the ultimate
  release law $(X_{\rm rel},Y_{\rm rel})=(X_{\tau_c-},Y_{\tau_c-})$ with
  joint pmf $\PR_i\{X_{\rm rel}=m,Y_{\rm rel}=n\}=\int_0^\infty(\delta_1m+
  \delta_2n)p^{(i)}_{m,n}(t)\dd t$ ($m+n\ge1$), atom $\sigma_i(\infty)$ at
  $(0,0)$.
\item \emph{Antal--Krapivsky recovery} (\texttt{app:AK-recovery}): at
  $\delta_1=\delta_2=0$ the pgf system reduces equation-by-equation to the
  Antal--Krapivsky (2011) two-type model under the rate map
  $\alpha^{\rm AK}_i=\lambda_i$, $\beta^{\rm AK}_i=\mu_i$,
  $\nu^{\rm AK}=\nu$.
\end{itemize}

\paragraph{Location and connections.} Chapter 7 (\texttt{tt:ch:two-type}).
Verified by Gillespie genealogies and RK4 integration against the closed
form. The one-type catastrophe limit collapses to M6's Riccati structure;
the GATE regime supplies one of the M20 skeleton instances.

% ---------------------------------------------------------------
\subsection{M29 --- Replicator--suppressor process}
\label{sec:M29}

\paragraph{Class.} Nonlinear CTMC (the interaction rate couples the two
coordinates; \emph{not} a branching process) on a lattice bounded in one
coordinate; a pure-birth process once the store is spent.

\paragraph{State space.} (\texttt{path:eq:statespace})
\[
\mathcal S=\bigl\{(i,j):\ i\in\mathbb N\cup\{0\},\ j\in\{0,\ldots,N\}\bigr\},
\]
$X_t$ = intracellular pathogens (replicators), $Q_t$ = remaining granules
(suppressor); $X_t=0$ absorbing (extinction), $Q_t=0$ absorbing for the
defence (pure birth thereafter).

\paragraph{Dynamics.} (\texttt{path:eq:rsrates}):
\[
\PR{\Delta X=1,\Delta Q=0}= \lambda i\,\Delta t \ \text{(birth)},\qquad
\PR{\Delta X=-1,\Delta Q=-1}= \varsigma ij\,\Delta t \ \text{(suppression)}
\]
plus the no-event term $1-(\lambda i+\varsigma ij)\Delta t$. Suppression
removes one replicator \emph{and} one granule --- ``that product makes the
process nonlinear, and it is the only nonlinearity in the model''.

\paragraph{Parameters.} $\lambda$ per-replicator birth; $\varsigma$
suppression rate per replicator--granule pair; $N$ granule store of a fresh
cell (no replenishment, Remark \texttt{path:rem:noreplenish}).

\paragraph{Initial conditions.} $X_0=i$ (incident dose), $Q_0=j$
(typically $j=N$).

\paragraph{Principal results.} Extinction probabilities
$\pi_{i,j}=\PR(\text{fixation in }X=0\mid(i,j))$ satisfy the embedded
first-step recursion (\texttt{path:eq:firststep})
\[
\pi_{i,j}=\hat\alpha_j\,\pi_{i+1,j}+\hat\beta_j\,\pi_{i-1,j-1},\qquad
\hat\alpha_j=\frac{\lambda}{\lambda+\varsigma j},\qquad
\hat\beta_j=\frac{\varsigma j}{\lambda+\varsigma j},
\]
with $\pi_{0,j}=1$, $\pi_{i,0}=0$ ($i\ge1$). \emph{Certainty threshold}
(Proposition \texttt{path:prop:certainty}): $\pi_{i,j}=0$ whenever $i>j$
--- survival is certain once the cohort exceeds the store. \emph{Diagonal
law} (Proposition \texttt{path:prop:diagonal}), with $\nu=\lambda/\varsigma$:
\[
\pi_{j,j}=\prod_{m=1}^{j}\frac{\varsigma m}{\lambda+\varsigma m}
=\frac{j!\,\Gamma(\nu+1)}{\Gamma(j+\nu+1)}
=\binom{j+\nu}{j}^{-1}\ \sim\ \Gamma(\nu+1)\,j^{-\nu},
\]
exactly $1/(j+1)$ at $\lambda=\varsigma$; off-diagonal $\pi_{i,j}$ has no
product form (open). Verified against a $210$-unknown linear solve at
$N=20$ and 400{,}000-path Gillespie of the embedded chain.

\paragraph{Location and connections.} Chapter 8
(\texttt{path:def:rs}). Intracellular counterpart of the flooding
dose-dependence (M23); pairs with M30/M31. The thesis notes the constant
per-granule-rate idealisation (death rate falls as the store is spent) as
the weakest link in the Chapter 10 quadrant argument (M38).

% ---------------------------------------------------------------
\subsection{M30 --- Extracellular ``playing field'' particle model}

\paragraph{Class.} Spatial stochastic individual-based simulation
(Brownian motion $+$ contact interactions); not formalised as an
equation-defined Markov process.

\paragraph{State space.} Three agent types in a 3-D extracellular matrix:
replicative units, target cells (uniform density), killer cells (denser,
smaller).

\paragraph{Dynamics.} Brownian movement; a killer cell that contacts a
pathogen eliminates it one-to-one and \emph{is used up}; inside target
cells the BDC process M6 runs until rupture releases the load. Quadtree
collision indexing (Appendix \texttt{path:app:quadtree}).

\paragraph{Parameters.} Spatial densities; burst size $r$; casualties
$\vartheta(r)$; killer prevalence assumed reset between ruptures (Remark
\texttt{path:rem:reset}), making $\vartheta$ iid across deliveries.

\paragraph{Principal results.} Saturating casualty law and its
idealisation (\texttt{path:eq:logremoval}, \texttt{path:eq:fixedremoval}):
\[
\vartheta(r)\sim\log(r+1),\qquad
\text{removal}=\min(r,\vartheta)=\begin{cases}r,&r\le\vartheta,\\
\vartheta,&r>\vartheta.\end{cases}
\]

\paragraph{Location and connections.} Chapter 8
\S\texttt{path:sec:playingfield}. Justifies the fixed-removal caricature
applied in M31.

% ---------------------------------------------------------------
\subsection{M31 --- Fixed-removal replication-ratio model}

\paragraph{Class.} Output functional over the M6 burst law with threshold
removal (a ``one-delivery'' offspring model); plus the matched-mean budding
reversal comparison.

\paragraph{Dynamics.} The release $r$ of one founding bacterium's BDC
process: $\PR{r=i}=\frac{\delta}{\lambda}a^{-i}$ ($i\ge1$), $\PR{r=0}=b$
(\texttt{path:eq:burstlaw}); survivors pass through
$r\mapsto\max(0,r-\vartheta)$.

\paragraph{Parameters.} $\lambda,\mu,\delta$; roots $a,b$; expected yield
$V_\infty$; removal budget $\vartheta\in\mathbb Z_{\ge0}$; threshold budget
$\vartheta^*=\log V_\infty/\log a$; flooding parameter
$L=a(1-b)$; budding match $\sigma=V_\infty/(1+V_\infty)$.

\paragraph{Principal results.} Exact upper tail
$\PR{r\ge\vartheta}=(1-b)a^{-(\vartheta-1)}$ (\texttt{path:eq:tail});
\emph{replication ratio} (Proposition \texttt{path:prop:ratio},
\texttt{path:eq:R0}) --- by memorylessness of the geometric tail,
\[
\mathcal R(\vartheta)=\EX{\max(0,r-\vartheta)}=V_\infty\,a^{-\vartheta},
\]
with $\mathcal R>1$ iff $\vartheta<\vartheta^*=\log V_\infty/\log a$
(Corollary \texttt{path:cor:threshold}); $\mathcal R(0)=V_\infty>1$ always,
so subcriticality exists only because of removal (Corollary
\texttt{path:cor:whythreshold}); real budgets:
$\EX{\max(0,r-\vartheta)}=a^{-\lceil\vartheta\rceil}\{V_\infty+
(\lceil\vartheta\rceil-\vartheta)L\}$. \emph{Reversal} (Proposition
\texttt{path:prop:reversal}, \texttt{path:eq:reversal}):
\[
\mathcal R_{\rm burst}(\vartheta)=V_\infty a^{-\vartheta},\qquad
\mathcal R_{\rm bud}(\vartheta)=V_\infty\sigma^{\vartheta},
\]
and bursting beats matched-mean budding iff $1/a>\sigma$ iff $L<1$ --- the
exact complement of the flooding theorem M23. Batch-count argument: a
yield delivered in $n$ instalments retains $\max(0,V_\infty-n\vartheta)$
against a burster's $\max(0,V_\infty-\vartheta)$ --- concentration wins for
every $n,\vartheta,L$.

\paragraph{Location and connections.} Chapter 8
\S\texttt{path:sec:ratio}, \S\texttt{path:sec:reversal}; verified against a
3000-state embedded-chain solve and Monte Carlo. Bridges M6/M16 to the
host-level switch narrative; the Chapter 10 open problem asks to reconcile
$\vartheta^*$ with M38's $p_N^*$.

% ============================================================================
\section{Catalogue, part III --- non-constant rates (Chapter 9) and conclusion (Chapter 10)}
% ============================================================================

% ---------------------------------------------------------------
\subsection{M32 --- Pimentel competition chain}

\paragraph{Class.} Discrete-time Markov chain with path-dependent kernel:
$A_t$ alone is not Markov; the triple $(A_t,R_{a,t},R_{b,t})$ is.

\paragraph{State space.} $A_t\in\{0,\ldots,N\}$ with $A_t+B_t=N$; $0$ and
$N$ absorbing; resistances $R_{a,t},R_{b,t}\in\mathbb R_+$.

\paragraph{Dynamics.} Kernel (\texttt{var:eq:kernel}):
\[
\PR{A_{t+1}=j\mid A_t=i}=
\begin{cases}
q_a=\frac iN(1-p_b), & j=i+1,\\[2pt]
q_b=(1-\frac iN)(1-p_a), & j=i-1,\\[2pt]
q_c=\frac iN p_b+(1-\frac iN)p_a, & j=i,
\end{cases}
\]
with blocking probabilities and resistance updates
(\texttt{var:eq:blockprobs}, \texttt{var:eq:resistupdate}):
\[
p_a=\max\{0,1-\tfrac{R_b}{R_a}\},\quad p_b=\max\{0,1-\tfrac{R_a}{R_b}\},
\quad
R_{k,t+1}=R_{k,t}+\varepsilon\,\Delta R_{k,t},
\]
\[
\Delta R_{a,t}=\Bigl(\frac{N-A_t}{A_t}\Bigr)^{\alpha},\qquad
\Delta R_{b,t}=\Bigl(\frac{A_t}{N-A_t}\Bigr)^{\alpha}.
\]

\paragraph{Parameters.} $N$ enclosure capacity; $\varepsilon$ evolution
rate (only $\varepsilon/R_0$ matters); $\alpha$ minority-squeeze exponent
($\alpha=1$ physical; $\alpha>1$ makes criticality visible).

\paragraph{Initial conditions.} $A_0=B_0=N/2$; $R_{a,0}=R_{b,0}=1$.

\paragraph{Principal results.} The moving critical point (equation (9.21)
= \texttt{var:eq:icrit}):
\[
i_c=\begin{cases}
N\,\dfrac{1-p_a}{2-p_a}, & p_a>0,\ p_b=0,\\[6pt]
N\,\dfrac{1}{2-p_b}, & p_b>0,\ p_a=0,
\end{cases}
\]
unstable ($q_a>q_b$ iff $i>i_c$) and path-dependent: ``the resistance
update drags the critical point after the leading population''. Fixation
probabilities/times not computed (open); simulations show one or two
reversals then absorption at $\alpha=1$ and order-ten reversals with
growing swings at large $\alpha$. A spatial version (cellular automaton,
\S\texttt{var:sec:automaton}) inherits the same local rule.

\paragraph{Location and connections.} Chapter 9 \S3
(\texttt{var:sec:pimentel}), motivated by Pimentel's fly
competition-reversal experiment. Recast with a fitness coupling in M33.

% ---------------------------------------------------------------
\subsection{M33 --- Pimentel$+$ (coupling speciation to competition) \quad [specification]}

\paragraph{Class.} Discrete-time competition chain coupled to two
continuous-time linear birth--death species counts; specified, never
simulated.

\paragraph{Dynamics.} Fitness kernel (\texttt{var:eq:fitkernel}):
\[
\PR{\Delta A_t=1}=\frac{A_t}{N}\,\frac{f_a}{\bar f},\qquad
\PR{\Delta A_t=-1}=\Bigl(1-\frac{A_t}{N}\Bigr)\frac{f_b}{\bar f},
\]
fitness update with species-count penalty (\texttt{var:eq:fitpenalty}):
\[
f_a(t+\Delta t)=f_a(t)+\varepsilon\Bigl(\frac{B_t}{A_t}\Bigr)^{\alpha}
-\gamma S^{(a)}_t,\qquad
f_b \text{ symmetric},
\]
species counts as birth--death processes (\texttt{var:eq:subspec}):
\[
\PR{\Delta S^{(a)}_t=\pm1}=\beta_aS^{(a)}_t\Delta t,\ \ \mu_aS^{(a)}_t
\Delta t,\qquad
\beta_a=\Bigl(\frac{A_t}{N}\Bigr)^{\!\zeta}f_a .
\]

\paragraph{Parameters.} $\varepsilon,\alpha$ as M32; $\mu_a$ per-species
extinction; $\zeta$ dominance exponent; $\gamma$ cost of carrying excess
variation.

\paragraph{Role.} Contains both reversal routes at once (minority squeeze
$+$ winner weakening); the $-\gamma S^{(a)}_t$ penalty is the unrun
ancestor of M35's per-type cost $-cX_i$. Open: whether the penalty term
alone can reverse; the regime where $\gamma S^{(a)}_t$ competes with
accumulated fitness.

\paragraph{Location.} Chapter 9 \S4 (\texttt{var:sec:pimentelplus}).

% ---------------------------------------------------------------
\subsection{M34 --- Rise, run, ruin, rejuvenation (RRRR)}

\paragraph{Class.} CTMC with state-dependent catastrophe rate plus an
auxiliary ODE --- the pair $(X_t,V(t))$ is Markov, $X_t$ alone is not.
The thesis's own classification: ``a birth--death--catastrophe process with
catastrophe rate $\delta X_t/V(t)$ and an auxiliary ODE for $V$''.

\paragraph{State space.} $X_t\in\mathbb N\cup\{0\}$ (replicative
agencies); $V(t)\in(0,\infty)$ accumulated potential.

\paragraph{Dynamics.} (\texttt{var:eq:rrrr}, \texttt{var:eq:rrrrV}):
\[
\PR{\Delta X_t=1}=\lambda X_t\Delta t \ \text{(proliferation)},\qquad
\PR{X_{t+\Delta t}=0}=\delta\,\frac{X_t}{V(t)}\Delta t \ \text{(catastrophe)},
\]
\[
\PR{X_{t+\Delta t}=1\mid X_t=0}=\omega\Delta t \ \text{(revival)},\qquad
\dd V/\dd t=\phi-cX_t .
\]
Revival keeps $0$ from being absorbing and makes the process recurrent; as
$V\downarrow0$ the hazard behaves like $\delta/\{c(t^*-t)\}$, which
integrates to infinity, so catastrophe arrives almost surely first.

\paragraph{Parameters.} $\phi$ potential inflow; $c$ consumption per
agency; $\lambda$, $\delta$, $\omega$ (simulation:
$(\phi,c,\lambda,\delta,\omega)=(4,0.8,0.5,0.25,0.15)$).

\paragraph{Principal results.} Simulation only: a four-movement cycle
(build, grow, consume, collapse) whose amplitude is fixed by the parameters
but whose period is random and de-phases across realisations. Open: the
cycle-period distribution (a first-passage problem for the pair), a
quasi-stationary law under the scarcity hazard, identifiability ($V$ enters
only through the ratio).

\paragraph{Location and connections.} Chapter 9 \S5
(\texttt{var:sec:rrrr}). M6 with the constant hazard $\delta X_t$ replaced
by $\delta X_t/V(t)$; first half of a matched pair with M35.

% ---------------------------------------------------------------
\subsection{M35 --- Multiplicative dissipation automaton}

\paragraph{Class.} Stochastic cellular automaton / interacting particle
system on a torus; ``the chapter's least tractable model''.

\paragraph{State space.} $L\times L$ lattice ($L=256$), Moore
neighbourhood, periodic boundary; each cell empty or occupied, occupants
carry a species and nested sub-species label; per extant species $i$:
potential $V_i$, cells held $N_i$, extant sub-species $X_i\le N_i$. Time in
sweeps.

\paragraph{Dynamics.} Contest strength of a cell $\zeta$ (sub-species $u$
of species $i$, \texttt{var:eq:castrength}):
\[
w(\zeta)=V_i\,(n_\zeta+1)^{\theta}\,\ee^{-\nu|u|/\kappa_i},
\qquad \kappa_i=\frac{N_i}{X_i},
\]
won with probability $w(\zeta)/(w(\zeta)+w(\zeta'))$; empty ground resists
with fixed strength $d$ (occupied cell takes a contested empty neighbour
w.p.\ $1-d$; occupied cells never vacate). Mutation rate per species
$R^{\rm mut}_i=mV_iN_i$ (\texttt{var:eq:camut}); potential ODE
(\texttt{var:eq:caode}) $\dd V_i/\dd t=\varphi N_i-cX_i$ floored at zero;
fragmentation hazard (\texttt{var:eq:cafrag}) $R^{\rm frag}_i=\delta X_i/V_i$
(plus a deterministic trigger when the ODE update would cross zero), at
which every sub-species becomes an independent species (potential reset to
$V_0$).

\paragraph{Parameters.} $L$, $V_0$, $\varphi$, $c$, $m$, $\theta$, $\nu$,
$d$, $\delta$.

\paragraph{Initial conditions.} One occupied central cell, one species, one
sub-species, potential $V_0$.

\paragraph{Principal results.} Balance locus (\texttt{var:eq:cabalance})
$X_i\approx\varphi N_i/c$ where production and dissipation cancel; a
species can lose potential only if $c>\varphi$; a single-sub-species
species needs $\ge c/\varphi$ cells to break even (burst filter by area).
Key structural observation: on the balance locus the hazard becomes
$\delta X_i/V_i\approx(\delta\varphi/c\bar V)N_i$ --- proportional to area
--- so the species count is \emph{conjectured} to read as M6's process with
burst size $X_i$ and effective $\delta$; supported by the founder burst at
$t=354$ (8689 cells, 964 sub-species) but not derived.

\paragraph{Location and connections.} Chapter 9 \S6
(\texttt{var:sec:ca}). Deliberate pair to M34 (inflow per cell and
endogenous; catastrophe divides rather than destroys); realises and
simulates the penalty mechanism M33 only specified.

% ---------------------------------------------------------------
\subsection{M36 --- Birth--plague model (two variants)}

\paragraph{Class.} Two-type CTMC with a bilinear infection term
(density-dependent rates); variant 2 couples an auxiliary vitality ODE into
the loss rate.

\paragraph{State space.} $(X_t,Y_t)\in\mathbb N^2$: uninfected and
infected individuals; variant adds $V(t)\in\mathbb R$.

\paragraph{Dynamics.} (\texttt{var:eq:plague}):
\[
\PR{\Delta X=1}=\lambda X_t\Delta t,\quad
\PR{\Delta X=-1,\Delta Y=1}=(\delta+\chi Y_t)X_t\Delta t,\quad
\PR{\Delta Y=-1}=\mu Y_t\Delta t,
\]
the middle channel superposing spontaneous corruption (rate $\delta$ per
head) and contact infection (rate $\chi$ per pair). Vitality variant
(\texttt{var:eq:vitality}, \texttt{var:eq:vitalityV}): other channels
unchanged, plus $\dd V/\dd t=\phi-cX_t$ and corruption slowed to
$\delta X_t/V(t)$ (by consistency with M34; the source flags an
alternative $\delta X_t^2/V(t)$ reading as ``author to confirm'').

\paragraph{Parameters.} $\lambda,\delta,\chi,\mu$ (simulation
$(1,0.1,0.02,10)$); variant $\phi,c$.

\paragraph{Principal results.} Base model simulated only (infection acts
as a brake, no collapse mechanism); the vitality variant is specified to
supply the missing amplification and has not been simulated.

\paragraph{Location and connections.} Chapter 9, Appendix 9.A
(\texttt{var:app:plague}). Complement to M32's coupling idea (loss rate
driven by a second count rather than a carried potential).

% ---------------------------------------------------------------
\subsection{M37 --- Public good accumulation}

\paragraph{Class.} CTMC birth--death process whose birth rate is a
functional of an auxiliary ODE compound (positive feedback); embedded in a
host-to-host Galton--Watson layer.

\paragraph{State space.} $X_t\in\mathbb N$ (bacteria within one
macrophage); $P(t)\in\mathbb R$ (public-good concentration).

\paragraph{Dynamics.} (\texttt{var:eq:publicgood}):
\[
\dd P/\dd t=gX_t-c,\qquad
\PR{\Delta X_t=1}=\lambda(t)X_t\Delta t,\quad
\lambda(t)=\lambda_0+\alpha P(t),\qquad
\PR{\Delta X_t=-1}=\mu X_t\Delta t .
\]

\paragraph{Parameters.} $g$ production per bacterium; $c$ clearance;
$\lambda_0$ base birth; $\alpha$ public-good advantage; $\mu$ death;
regime of interest $\mu>\lambda_0$ (subcritical alone, supercritical once
enough good accumulates; while $X_t<c/g$ the good is cleared faster than
made).

\paragraph{Principal results.} Minimum founding cohort
(\texttt{var:eq:kstar}):
\[
K^*=\inf\Bigl\{K:\ \lim_{t\to\infty}\PR{X_t>0\mid X_0=K}>\tfrac12\Bigr\},
\]
computed numerically (steep sigmoid; $K^*$ rises with $\mu-\lambda_0$ and
sits above $c/g$; at $\mu-\lambda_0=1.2$ no $K\le28$ crossed one half).
Host-to-host embedding: cells that survive to late time deliver their load
into two more cells --- a Galton--Watson process (M2 with general
offspring law), supercritical exactly when within-cell survival exceeds
one half.

\paragraph{Location and connections.} Chapter 9, Appendix 9.B
(\texttt{var:app:publicgood}). The feedback run in the favourable
direction (complement of M34/M36); founding-cohort emphasis matches the
thesis-wide near-criticality theme.

% ---------------------------------------------------------------
\subsection{M38 --- Quadrant argument \quad [D: comparison]}

\paragraph{Class.} \textbf{[D]} Comparison of expected net drifts of three
simple birth--death processes; ``the crudest argument in the thesis''.

\paragraph{Setting.} Environments $i\in\{M,N,E\}$ (macrophage interior,
neutrophil interior, extracellular), each a linear birth--death process
with rates $(\beta_i,\mu_i)$; orderings
\[
\mu_N>\mu_E>\mu_M,\qquad \beta_M=\beta_E=\beta_N=\beta,\qquad
\beta_M-\mu_M>\beta_E-\mu_E>\beta_N-\mu_N .
\]

\paragraph{Dynamics of the comparison.} With phagocyte composition
$p_M=\#M/(\#M+\#N)$, $p_N=1-p_M$: a pro-phagocytic strategist draws its
environment at random from the two interiors, an anti-phagocytic one takes
the extracellular rate with certainty:
\[
R_{\rm pro}=p_M(\beta_M-\mu_M)+p_N(\beta_N-\mu_N),\qquad
R_{\rm anti}=\beta_E-\mu_E .
\]

\paragraph{Principal result.} Pro-phagocytosis has the larger expected
drift iff (boxed, \texttt{con:eq:threshold})
\[
p_N>p_N^*:=\frac{\mu_E-\mu_M}{\mu_N-\mu_M},
\]
the position of the extracellular death rate on the interval from the
macrophage to the neutrophil rate; the switch is keyed to composition, not
elapsed time, and is testable by neutrophil depletion/recruitment.

\paragraph{Location and connections.} Chapter 10
\S\texttt{con:sec:quadrant}. Companion threshold to M31's
$\vartheta^*=\log V_\infty/\log a$ --- ``reached from a different model
with different machinery, and the two accounts have not been reconciled''
(Chapter 10 open problem 1). Weakest stated link: it treats $\mu_M,\mu_N$
as constants where M29 shows they deplete with the granule store.

% ============================================================================
\section{Solution methods the thesis builds around}
% ============================================================================

Not models, but the toolkit the entries above are solved with; recorded
here for completeness.

\begin{itemize}
\item \emph{First-step principle} (\texttt{m:prop:firststep}):
  $h_i=\sum_jP_{ij}h_j$ (discrete) and $h_i=\sum_{j\ne i}\frac{q_{ij}}
  {q_i}h_j$ (continuous), with absorption-time variants; branching fixed
  point $s=\phi(s)$ (\texttt{m:eq:extinctionfixedpoint}). Used by M1, M5,
  M23, M29, M32.
\item \emph{Probability generating functions}: composition for random sums;
  bivariate pgfs (\texttt{m:eq:pgfbivariate}); coefficient extraction and
  the bivariate Cauchy integral (Appendix 2.D). Used by M2, M5, M6,
  M13--M16, M28.
\item \emph{Method of characteristics} (\texttt{m:sec:moc}): generic
  first-order PDE machinery anchoring to an initial surface; the worked
  example is M14, extended to the Riccati characteristic of M15 and the
  pgf PDE of M10.
\item \emph{Riccati root factorisation}: the quadratic
  $\lambda f^2-(\lambda+\mu+\delta)f+\mu$ with roots $a,b$ (M6); the same
  structure recurs in M8, M15, M28 (hypergeometric and Bessel reductions).
\item \emph{Feynman--Kac / occupation-time transforms} (M28), defective
  (killed) generating functions (M6, M28), time-change of the Yule law
  (M4 $\to$ M10).
\item \emph{Simulation}: Gillespie's direct method
  (\texttt{m:rem:gillespie}) with competing exponential clocks; Poisson
  thinning for time-dependent rates (M9); embedded-chain truncation solves
  and Monte Carlo as verification (Chapters 5, 6, 8).
\end{itemize}

% ============================================================================
\section{How the models connect}
% ============================================================================

\begin{itemize}
\item \textbf{The BDC spine.} M6 (Chapters 2, 4, 5) is the thesis's
  engine: its roots $a,b$ ``carry every closed form''; its kernels
  $\hat I$ and $g=\delta K$ drive the population layer M20; its burst law
  feeds the establishment process M23 and the removal functional M31; its
  two-type extension is M28; its scarcity-rate deformation is M34 (and,
  conjecturally, M35).
\item \textbf{Discrete--continuous pair.} M2 (Galton--Watson, $A(p)$) and
  M5 (birth--death, $A_c$) pose the same quasi-stationary question in
  discrete and continuous time; $A_c$ is elementary, $A(p)$ is not, and the
  Koenigs function explains why (M2).
\item \textbf{Budding vs bursting.} M17 vs M6 at the cell level; M23
  (flooding, $L>1$ favours bursting without a budget) and M31 (reversal,
  $L<1$ favours bursting under threshold removal) are exact complements;
  M20--M22 show the schedule is invisible to $R_0$ and to mean dynamics.
\item \textbf{Absorption family.} M13--M15 model post-rupture transfer
  into fresh compartments; M15's closed form is a thesis result; M16 chains
  the cells instead.
\item \textbf{Non-constant rates.} Chapter 9's taxonomy: what costs the
  generating function is ``continuous bending of a rate''. M10 (moving
  birth, exact geometric law), M32/M33 (path-dependent kernels), M34--M37
  (feedback through auxiliary compounds or second counts). The conclusion
  restates the principle: a process may act on its own circumstances and
  stay solvable, provided what survives the interaction is a random time
  or a random count rather than a moving rate.
\item \textbf{Thresholds.} Three independent thresholds appear:
  $\vartheta^*=\log V_\infty/\log a$ (M31), $p_N^*$ (M38), $K^*$
  (M37); reconciling the first two is an open problem of the conclusion.
\end{itemize}

\vfill
\begin{center}\small
Extracted September 2026 from \texttt{full\_draft\_thesis} (stitched working
draft). Thesis \texttt{label} keys are quoted throughout as stable
identifiers; all equations are reproduced from the chapter sources.
\end{center}

\end{document}
